%%%%%%%%%%%%%%%%%%%%%%%%%%%%%%%%%%%%%%%%%
% Journal Article
% LaTeX Template
% Version 1.4 (15/5/16)
%
% This template has been downloaded from:
% http://www.LaTeXTemplates.com
%
% Original author:
% Frits Wenneker (http://www.howtotex.com) with extensive modifications by
% Vel (vel@LaTeXTemplates.com)
%
% License:
% CC BY-NC-SA 3.0 (http://creativecommons.org/licenses/by-nc-sa/3.0/)
%
%%%%%%%%%%%%%%%%%%%%%%%%%%%%%%%%%%%%%%%%%

%----------------------------------------------------------------------------------------
%	PACKAGES AND OTHER DOCUMENT CONFIGURATIONS
%----------------------------------------------------------------------------------------

\documentclass[twoside,twocolumn]{article}

\usepackage{blindtext} % Package to generate dummy text throughout this template 

\usepackage{pgf}
\usepackage[sc]{mathpazo} % Use the Palatino font
\usepackage{longtable}
\usepackage[T1]{fontenc} % Use 8-bit encoding that has 256 glyphs
\linespread{1.05} % Line spacing - Palatino needs more space between lines
\usepackage{microtype} % Slightly tweak font spacing for aesthetics

\usepackage[english]{babel} % Language hyphenation and typographical rules

\usepackage[hmarginratio=1:1,top=32mm,columnsep=20pt]{geometry} % Document margins
\usepackage[hang, small,labelfont=bf,up,textfont=it,up]{caption} % Custom captions under/above floats in tables or figures
\usepackage{booktabs} % Horizontal rules in tables

\usepackage{lettrine} % The lettrine is the first enlarged letter at the beginning of the text
\usepackage{amsmath}
\usepackage{enumitem} % Customized lists
\setlist[itemize]{noitemsep} % Make itemize lists more compact

\usepackage{abstract} % Allows abstract customization
\renewcommand{\abstractnamefont}{\normalfont\bfseries} % Set the "Abstract" text to bold
\renewcommand{\abstracttextfont}{\normalfont\small\itshape} % Set the abstract itself to small italic text

\usepackage{titlesec} % Allows customization of titles
\renewcommand\thesection{\Roman{section}} % Roman numerals for the sections
\renewcommand\thesubsection{\roman{subsection}} % roman numerals for subsections
\titleformat{\section}[block]{\large\scshape\centering}{\thesection.}{1em}{} % Change the look of the section titles
\titleformat{\subsection}[block]{\normalsize\itshape}{\thesubsection. }{0em}{} % Change the look of the section titles
\titlespacing*{\subsection}{0pt}{2pt}{1pt}

\usepackage{fancyhdr} % Headers and footers
\pagestyle{fancy} % All pages have headers and footers
\fancyhead{} % Blank out the default header
\fancyfoot{} % Blank out the default footer
\fancyhead[C]{Churn Prediction $\bullet$ December 2021 $\bullet$ Introduction to DS 1001} % Custom header text
\fancyfoot[RO,LE]{\thepage} % Custom footer text

\usepackage{titling} % Customizing the title section
\usepackage{layouts}
\usepackage{hyperref} % For hyperlinks in the PDF

%----------------------------------------------------------------------------------------
%	TITLE SECTION
%----------------------------------------------------------------------------------------

\setlength{\droptitle}{-4\baselineskip} % Move the title up

\pretitle{\begin{center}\Huge\bfseries} % Article title formatting
\posttitle{\end{center}} % Article title closing formatting
\title{Churn Prediction in Retail Banking Customers} % Article title
\author{%
\textsc{Adi Srikanth $\bullet$ Andre Chen $\bullet$ David Roth} \\[1ex] % Your name
\textsc{Jin Ishizuka $\bullet$ Lev Paciorkowski} \\[1ex] % Your name
\normalsize NYU Center for Data Science \\ % Your institution
% \normalsize \href{mailto:john@smith.com}{john@smith.com} % Your email address
%\and % Uncomment if 2 authors are required, duplicate these 4 lines if more
%\textsc{Jane Smith}\thanks{Corresponding author} \\[1ex] % Second author's name
%\normalsize University of Utah \\ % Second author's institution
%\normalsize \href{mailto:jane@smith.com}{jane@smith.com} % Second author's email address
}
\date{\today} % Leave empty to omit a date
\renewcommand{\maketitlehookd}{%
\begin{abstract}
\noindent Customer "churn" refers to the rate at which...
\end{abstract}
}

%----------------------------------------------------------------------------------------

\begin{document}

% Print the title
\maketitle

%----------------------------------------------------------------------------------------
%	ARTICLE CONTENTS
%----------------------------------------------------------------------------------------
\section{Introduction}

\lettrine[nindent=0em,lines=2]{I} n this analysis, we look at a dataset of $10,000$ retail banking customers and try to predict whether an individual customer is likely to "churn", that is, if they are likely to close their account and cease their financial relationship with their bank.

\section{Data Summary}

\lettrine[nindent=0em,lines=2]{O}UR dataset contains information on customer churn at a retail bank, specifically features/predictor variables and a target variable ($0$ if the customer churned, $1$ if the customer is still at the bank). The features include general information about the customers including demographics, socioeconomic information, and the level of engagement each customer has with the bank. Our goal is to study the relationships between the customer predictor variables and their probability of churning to understand what leads to customer churn at this bank.
    
We determined that investigating this dataset would be broadly useful as practice for future roles as data scientists in industry. Specifically, a very common use case in both consumer and enterprise-facing companies is understanding leading indicators of churn, as churn contributes to significant revenue erosion and needs to be proactively addressed. This dataset offers a great opportunity for us to practice not only the technical skills involved in determining predictors of customer churn, but also the more business-oriented aspects of asking suitable questions and presenting a cohesive conclusion that could benefit our hypothetical company.

\noindent The data is broken down as follows:\smallskip\\
\noindent\textbf{Rows}: Each row represents an individual customer of the bank, providing information about that customer and whether or not they churned, which we assume in this context means they closed their account with the bank and stopped doing business with them altogether.\smallskip\\
\noindent\textbf{Columns}: The 14th and last column of the dataset (“Exited”) is a binary target variable that indicates whether a customer churned (0) or not (1). The first 13 columns of the dataset include information about each individual customer, of which a large subset were identified as potentially useful predictor variables and are described in detail below. We omitted columns like “CustomerId” and “Surname” from the columns listed below, as these columns likely have no relationship with a customer’s likelihood of churning:
\begin{itemize}
  \item \textit{CreditScore}: The customer’s credit score on a scale from 350 (worst) to 850 (best)
  \item \textit{Geography}: The customer’s country of residence (dataset includes France, Germany, and Spain)
  \item \textit{Gender}: Male or Female
  \item \textit{Age}: Age of customer $(18-92)$
  \item \textit{Tenure}: How long the customer has been with the bank, in years.
  \item \textit{Balance}: Customer balance to date. (Note: customers who have closed their accounts with the bnk may still have uncollected balances)
  \item \textit{NumOfProducts}: The number of bank products used by the customer.
  \item \textit{HasCrCard}: Whether the customer has a credit card $(1)$ or not $(0)$
  \item \textit{IsActiveMember}: Binary variable indicating if customer's activity is above the bank's minimim threshold for ``Active''.
  \item \textit{Estimated Salary}: Customer's estimated yearly salary.
\end{itemize}

\section{Hypothesis Testing}
\small \textit{Do customers with varying levels of engagement with the bank have significantly different probabilities of churn? What does this tell us about the measures of engagement which matter the most?}\medskip

\lettrine[nindent=0em,lines=2]{O}UR ultimate goal is to determine indicators of churn that the bank can act on, so that we could then recommend specific areas the bank could target for improvement in their interactions with customers. We accomplish this by defining customer profiles based on three measures of engagement (\textit{NumOfProducts}, \textit{IsActiveMember} and \textit{Tenure}), and then implement hypothesis testing to see if churn proportion is significantly different from one profile to another.

\subsection*{``Product-Driven'' Customers}
We find that only about 3\% of customers have more than two products with the bank, and of them, 85.9\% churned compared to just 18.2\% of non Product-Driven customers.
\begin{align*}
  H_0: P(churn)_{\text{\tiny{product-driven}}} &= P(churn)_{\text{\tiny{non-product-driven}}}\\
  H_1: P(churn)_{\text{\tiny{product-driven}}} &> P(churn)_{\text{\tiny{non-product-driven}}}
\end{align*}

\subsection*{``Active'' Customers}
Customers are evenly split between Active and not Active (as determined by the bank), and only 14.3\% of Active customers churned, compared to 26.9\% of non-Active customers.
\begin{align*}
  H_0: P(churn)_{\text{\tiny{active}}} &= P(churn)_{\text{\tiny{non-active}}}\\
  H_1: P(churn)_{\text{\tiny{active}}} &> P(churn)_{\text{\tiny{non-active}}}
\end{align*}

\subsection*{``Established'' Customers}
We find that Customers with tenures longer than 1 year churned at a rate of 20\%, compared to new customers, who churned at a rate of 22.6\%.
\begin{align*}
  H_0: P(churn)_{\text{\tiny{established}}} &= P(churn)_{\text{\tiny{new}}}\\
  H_1: P(churn)_{\text{\tiny{established}}} &> P(churn)_{\text{\tiny{new}}}
\end{align*}

\subsection*{\large{Results}}
In general, Product-Driveness and “Active”-ness are highly significant predictors of churn at $\alpha = 0.005$, while Tenure is not (See Table \ref{table:significance}).  Although Product-Driven customers are quite rare, they have a much greater tendency to churn compared to mainstream customers who use only one or two products. Perhaps this could be an indication of so-called credit card “hopping”, but deeper investigation would be required to evaluate this. Based on these results, the bank can be confident in how it assesses “Active” customers, but might consider interventions to disincentivize customers from becoming Product-Driven. In contrast, a customer’s Tenure should not be a central focus.
\begin{table}
  \centering
  \caption{Significance Results}
  \begin{tabular}{llr}
    \toprule
    {} &                Test &        p-value \\
    \midrule
    0 &  \textit{Product-Drivenness} &  $2.9e-196$ \\
    1 &     \textit{''Active''-ness} &   $3.0e-55$ \\
    2 &              \textit{Tenure} &   $1.2e-02$ \\
    \bottomrule  
  \end{tabular}
  \label{table:significance}
\end{table}

\begin{figure}
  \begin{center}
    \resizebox{7cm}{!}{
      %% Creator: Matplotlib, PGF backend
%%
%% To include the figure in your LaTeX document, write
%%   \input{<filename>.pgf}
%%
%% Make sure the required packages are loaded in your preamble
%%   \usepackage{pgf}
%%
%% Also ensure that all the required font packages are loaded; for instance,
%% the lmodern package is sometimes necessary when using math font.
%%   \usepackage{lmodern}
%%
%% Figures using additional raster images can only be included by \input if
%% they are in the same directory as the main LaTeX file. For loading figures
%% from other directories you can use the `import` package
%%   \usepackage{import}
%%
%% and then include the figures with
%%   \import{<path to file>}{<filename>.pgf}
%%
%% Matplotlib used the following preamble
%%   \usepackage{fontspec}
%%
\begingroup%
\makeatletter%
\begin{pgfpicture}%
\pgfpathrectangle{\pgfpointorigin}{\pgfqpoint{6.000000in}{4.000000in}}%
\pgfusepath{use as bounding box, clip}%
\begin{pgfscope}%
\pgfsetbuttcap%
\pgfsetmiterjoin%
\pgfsetlinewidth{0.000000pt}%
\definecolor{currentstroke}{rgb}{1.000000,1.000000,1.000000}%
\pgfsetstrokecolor{currentstroke}%
\pgfsetstrokeopacity{0.000000}%
\pgfsetdash{}{0pt}%
\pgfpathmoveto{\pgfqpoint{0.000000in}{0.000000in}}%
\pgfpathlineto{\pgfqpoint{6.000000in}{0.000000in}}%
\pgfpathlineto{\pgfqpoint{6.000000in}{4.000000in}}%
\pgfpathlineto{\pgfqpoint{0.000000in}{4.000000in}}%
\pgfpathlineto{\pgfqpoint{0.000000in}{0.000000in}}%
\pgfpathclose%
\pgfusepath{}%
\end{pgfscope}%
\begin{pgfscope}%
\pgfsetbuttcap%
\pgfsetmiterjoin%
\definecolor{currentfill}{rgb}{1.000000,1.000000,1.000000}%
\pgfsetfillcolor{currentfill}%
\pgfsetlinewidth{0.000000pt}%
\definecolor{currentstroke}{rgb}{0.000000,0.000000,0.000000}%
\pgfsetstrokecolor{currentstroke}%
\pgfsetstrokeopacity{0.000000}%
\pgfsetdash{}{0pt}%
\pgfpathmoveto{\pgfqpoint{1.200000in}{1.200000in}}%
\pgfpathlineto{\pgfqpoint{4.560000in}{1.200000in}}%
\pgfpathlineto{\pgfqpoint{4.560000in}{3.520000in}}%
\pgfpathlineto{\pgfqpoint{1.200000in}{3.520000in}}%
\pgfpathlineto{\pgfqpoint{1.200000in}{1.200000in}}%
\pgfpathclose%
\pgfusepath{fill}%
\end{pgfscope}%
\begin{pgfscope}%
\pgfpathrectangle{\pgfqpoint{1.200000in}{1.200000in}}{\pgfqpoint{3.360000in}{2.320000in}}%
\pgfusepath{clip}%
\pgfsetbuttcap%
\pgfsetroundjoin%
\definecolor{currentfill}{rgb}{0.981377,0.920617,0.865369}%
\pgfsetfillcolor{currentfill}%
\pgfsetlinewidth{0.000000pt}%
\definecolor{currentstroke}{rgb}{1.000000,1.000000,1.000000}%
\pgfsetstrokecolor{currentstroke}%
\pgfsetdash{}{0pt}%
\pgfpathmoveto{\pgfqpoint{1.200000in}{3.520000in}}%
\pgfpathlineto{\pgfqpoint{1.536000in}{3.520000in}}%
\pgfpathlineto{\pgfqpoint{1.536000in}{3.288000in}}%
\pgfpathlineto{\pgfqpoint{1.200000in}{3.288000in}}%
\pgfpathlineto{\pgfqpoint{1.200000in}{3.520000in}}%
\pgfusepath{fill}%
\end{pgfscope}%
\begin{pgfscope}%
\pgfpathrectangle{\pgfqpoint{1.200000in}{1.200000in}}{\pgfqpoint{3.360000in}{2.320000in}}%
\pgfusepath{clip}%
\pgfsetbuttcap%
\pgfsetroundjoin%
\definecolor{currentfill}{rgb}{0.336898,0.117591,0.308894}%
\pgfsetfillcolor{currentfill}%
\pgfsetlinewidth{0.000000pt}%
\definecolor{currentstroke}{rgb}{1.000000,1.000000,1.000000}%
\pgfsetstrokecolor{currentstroke}%
\pgfsetdash{}{0pt}%
\pgfpathmoveto{\pgfqpoint{1.536000in}{3.520000in}}%
\pgfpathlineto{\pgfqpoint{1.872000in}{3.520000in}}%
\pgfpathlineto{\pgfqpoint{1.872000in}{3.288000in}}%
\pgfpathlineto{\pgfqpoint{1.536000in}{3.288000in}}%
\pgfpathlineto{\pgfqpoint{1.536000in}{3.520000in}}%
\pgfusepath{fill}%
\end{pgfscope}%
\begin{pgfscope}%
\pgfpathrectangle{\pgfqpoint{1.200000in}{1.200000in}}{\pgfqpoint{3.360000in}{2.320000in}}%
\pgfusepath{clip}%
\pgfsetbuttcap%
\pgfsetroundjoin%
\definecolor{currentfill}{rgb}{0.336898,0.117591,0.308894}%
\pgfsetfillcolor{currentfill}%
\pgfsetlinewidth{0.000000pt}%
\definecolor{currentstroke}{rgb}{1.000000,1.000000,1.000000}%
\pgfsetstrokecolor{currentstroke}%
\pgfsetdash{}{0pt}%
\pgfpathmoveto{\pgfqpoint{1.872000in}{3.520000in}}%
\pgfpathlineto{\pgfqpoint{2.208000in}{3.520000in}}%
\pgfpathlineto{\pgfqpoint{2.208000in}{3.288000in}}%
\pgfpathlineto{\pgfqpoint{1.872000in}{3.288000in}}%
\pgfpathlineto{\pgfqpoint{1.872000in}{3.520000in}}%
\pgfusepath{fill}%
\end{pgfscope}%
\begin{pgfscope}%
\pgfpathrectangle{\pgfqpoint{1.200000in}{1.200000in}}{\pgfqpoint{3.360000in}{2.320000in}}%
\pgfusepath{clip}%
\pgfsetbuttcap%
\pgfsetroundjoin%
\definecolor{currentfill}{rgb}{0.330670,0.117012,0.306326}%
\pgfsetfillcolor{currentfill}%
\pgfsetlinewidth{0.000000pt}%
\definecolor{currentstroke}{rgb}{1.000000,1.000000,1.000000}%
\pgfsetstrokecolor{currentstroke}%
\pgfsetdash{}{0pt}%
\pgfpathmoveto{\pgfqpoint{2.208000in}{3.520000in}}%
\pgfpathlineto{\pgfqpoint{2.544000in}{3.520000in}}%
\pgfpathlineto{\pgfqpoint{2.544000in}{3.288000in}}%
\pgfpathlineto{\pgfqpoint{2.208000in}{3.288000in}}%
\pgfpathlineto{\pgfqpoint{2.208000in}{3.520000in}}%
\pgfusepath{fill}%
\end{pgfscope}%
\begin{pgfscope}%
\pgfpathrectangle{\pgfqpoint{1.200000in}{1.200000in}}{\pgfqpoint{3.360000in}{2.320000in}}%
\pgfusepath{clip}%
\pgfsetbuttcap%
\pgfsetroundjoin%
\definecolor{currentfill}{rgb}{0.387280,0.120898,0.327118}%
\pgfsetfillcolor{currentfill}%
\pgfsetlinewidth{0.000000pt}%
\definecolor{currentstroke}{rgb}{1.000000,1.000000,1.000000}%
\pgfsetstrokecolor{currentstroke}%
\pgfsetdash{}{0pt}%
\pgfpathmoveto{\pgfqpoint{2.544000in}{3.520000in}}%
\pgfpathlineto{\pgfqpoint{2.880000in}{3.520000in}}%
\pgfpathlineto{\pgfqpoint{2.880000in}{3.288000in}}%
\pgfpathlineto{\pgfqpoint{2.544000in}{3.288000in}}%
\pgfpathlineto{\pgfqpoint{2.544000in}{3.520000in}}%
\pgfusepath{fill}%
\end{pgfscope}%
\begin{pgfscope}%
\pgfpathrectangle{\pgfqpoint{1.200000in}{1.200000in}}{\pgfqpoint{3.360000in}{2.320000in}}%
\pgfusepath{clip}%
\pgfsetbuttcap%
\pgfsetroundjoin%
\definecolor{currentfill}{rgb}{0.368272,0.119937,0.320767}%
\pgfsetfillcolor{currentfill}%
\pgfsetlinewidth{0.000000pt}%
\definecolor{currentstroke}{rgb}{1.000000,1.000000,1.000000}%
\pgfsetstrokecolor{currentstroke}%
\pgfsetdash{}{0pt}%
\pgfpathmoveto{\pgfqpoint{2.880000in}{3.520000in}}%
\pgfpathlineto{\pgfqpoint{3.216000in}{3.520000in}}%
\pgfpathlineto{\pgfqpoint{3.216000in}{3.288000in}}%
\pgfpathlineto{\pgfqpoint{2.880000in}{3.288000in}}%
\pgfpathlineto{\pgfqpoint{2.880000in}{3.520000in}}%
\pgfusepath{fill}%
\end{pgfscope}%
\begin{pgfscope}%
\pgfpathrectangle{\pgfqpoint{1.200000in}{1.200000in}}{\pgfqpoint{3.360000in}{2.320000in}}%
\pgfusepath{clip}%
\pgfsetbuttcap%
\pgfsetroundjoin%
\definecolor{currentfill}{rgb}{0.010608,0.018082,0.100187}%
\pgfsetfillcolor{currentfill}%
\pgfsetlinewidth{0.000000pt}%
\definecolor{currentstroke}{rgb}{1.000000,1.000000,1.000000}%
\pgfsetstrokecolor{currentstroke}%
\pgfsetdash{}{0pt}%
\pgfpathmoveto{\pgfqpoint{3.216000in}{3.520000in}}%
\pgfpathlineto{\pgfqpoint{3.552000in}{3.520000in}}%
\pgfpathlineto{\pgfqpoint{3.552000in}{3.288000in}}%
\pgfpathlineto{\pgfqpoint{3.216000in}{3.288000in}}%
\pgfpathlineto{\pgfqpoint{3.216000in}{3.520000in}}%
\pgfusepath{fill}%
\end{pgfscope}%
\begin{pgfscope}%
\pgfpathrectangle{\pgfqpoint{1.200000in}{1.200000in}}{\pgfqpoint{3.360000in}{2.320000in}}%
\pgfusepath{clip}%
\pgfsetbuttcap%
\pgfsetroundjoin%
\definecolor{currentfill}{rgb}{0.438608,0.121770,0.341195}%
\pgfsetfillcolor{currentfill}%
\pgfsetlinewidth{0.000000pt}%
\definecolor{currentstroke}{rgb}{1.000000,1.000000,1.000000}%
\pgfsetstrokecolor{currentstroke}%
\pgfsetdash{}{0pt}%
\pgfpathmoveto{\pgfqpoint{3.552000in}{3.520000in}}%
\pgfpathlineto{\pgfqpoint{3.888000in}{3.520000in}}%
\pgfpathlineto{\pgfqpoint{3.888000in}{3.288000in}}%
\pgfpathlineto{\pgfqpoint{3.552000in}{3.288000in}}%
\pgfpathlineto{\pgfqpoint{3.552000in}{3.520000in}}%
\pgfusepath{fill}%
\end{pgfscope}%
\begin{pgfscope}%
\pgfpathrectangle{\pgfqpoint{1.200000in}{1.200000in}}{\pgfqpoint{3.360000in}{2.320000in}}%
\pgfusepath{clip}%
\pgfsetbuttcap%
\pgfsetroundjoin%
\definecolor{currentfill}{rgb}{0.368272,0.119937,0.320767}%
\pgfsetfillcolor{currentfill}%
\pgfsetlinewidth{0.000000pt}%
\definecolor{currentstroke}{rgb}{1.000000,1.000000,1.000000}%
\pgfsetstrokecolor{currentstroke}%
\pgfsetdash{}{0pt}%
\pgfpathmoveto{\pgfqpoint{3.888000in}{3.520000in}}%
\pgfpathlineto{\pgfqpoint{4.224000in}{3.520000in}}%
\pgfpathlineto{\pgfqpoint{4.224000in}{3.288000in}}%
\pgfpathlineto{\pgfqpoint{3.888000in}{3.288000in}}%
\pgfpathlineto{\pgfqpoint{3.888000in}{3.520000in}}%
\pgfusepath{fill}%
\end{pgfscope}%
\begin{pgfscope}%
\pgfpathrectangle{\pgfqpoint{1.200000in}{1.200000in}}{\pgfqpoint{3.360000in}{2.320000in}}%
\pgfusepath{clip}%
\pgfsetbuttcap%
\pgfsetroundjoin%
\definecolor{currentfill}{rgb}{0.355676,0.119109,0.316210}%
\pgfsetfillcolor{currentfill}%
\pgfsetlinewidth{0.000000pt}%
\definecolor{currentstroke}{rgb}{1.000000,1.000000,1.000000}%
\pgfsetstrokecolor{currentstroke}%
\pgfsetdash{}{0pt}%
\pgfpathmoveto{\pgfqpoint{4.224000in}{3.520000in}}%
\pgfpathlineto{\pgfqpoint{4.560000in}{3.520000in}}%
\pgfpathlineto{\pgfqpoint{4.560000in}{3.288000in}}%
\pgfpathlineto{\pgfqpoint{4.224000in}{3.288000in}}%
\pgfpathlineto{\pgfqpoint{4.224000in}{3.520000in}}%
\pgfusepath{fill}%
\end{pgfscope}%
\begin{pgfscope}%
\pgfpathrectangle{\pgfqpoint{1.200000in}{1.200000in}}{\pgfqpoint{3.360000in}{2.320000in}}%
\pgfusepath{clip}%
\pgfsetbuttcap%
\pgfsetroundjoin%
\definecolor{currentfill}{rgb}{0.336898,0.117591,0.308894}%
\pgfsetfillcolor{currentfill}%
\pgfsetlinewidth{0.000000pt}%
\definecolor{currentstroke}{rgb}{1.000000,1.000000,1.000000}%
\pgfsetstrokecolor{currentstroke}%
\pgfsetdash{}{0pt}%
\pgfpathmoveto{\pgfqpoint{1.200000in}{3.288000in}}%
\pgfpathlineto{\pgfqpoint{1.536000in}{3.288000in}}%
\pgfpathlineto{\pgfqpoint{1.536000in}{3.056000in}}%
\pgfpathlineto{\pgfqpoint{1.200000in}{3.056000in}}%
\pgfpathlineto{\pgfqpoint{1.200000in}{3.288000in}}%
\pgfusepath{fill}%
\end{pgfscope}%
\begin{pgfscope}%
\pgfpathrectangle{\pgfqpoint{1.200000in}{1.200000in}}{\pgfqpoint{3.360000in}{2.320000in}}%
\pgfusepath{clip}%
\pgfsetbuttcap%
\pgfsetroundjoin%
\definecolor{currentfill}{rgb}{0.981377,0.920617,0.865369}%
\pgfsetfillcolor{currentfill}%
\pgfsetlinewidth{0.000000pt}%
\definecolor{currentstroke}{rgb}{1.000000,1.000000,1.000000}%
\pgfsetstrokecolor{currentstroke}%
\pgfsetdash{}{0pt}%
\pgfpathmoveto{\pgfqpoint{1.536000in}{3.288000in}}%
\pgfpathlineto{\pgfqpoint{1.872000in}{3.288000in}}%
\pgfpathlineto{\pgfqpoint{1.872000in}{3.056000in}}%
\pgfpathlineto{\pgfqpoint{1.536000in}{3.056000in}}%
\pgfpathlineto{\pgfqpoint{1.536000in}{3.288000in}}%
\pgfusepath{fill}%
\end{pgfscope}%
\begin{pgfscope}%
\pgfpathrectangle{\pgfqpoint{1.200000in}{1.200000in}}{\pgfqpoint{3.360000in}{2.320000in}}%
\pgfusepath{clip}%
\pgfsetbuttcap%
\pgfsetroundjoin%
\definecolor{currentfill}{rgb}{0.318263,0.115743,0.301001}%
\pgfsetfillcolor{currentfill}%
\pgfsetlinewidth{0.000000pt}%
\definecolor{currentstroke}{rgb}{1.000000,1.000000,1.000000}%
\pgfsetstrokecolor{currentstroke}%
\pgfsetdash{}{0pt}%
\pgfpathmoveto{\pgfqpoint{1.872000in}{3.288000in}}%
\pgfpathlineto{\pgfqpoint{2.208000in}{3.288000in}}%
\pgfpathlineto{\pgfqpoint{2.208000in}{3.056000in}}%
\pgfpathlineto{\pgfqpoint{1.872000in}{3.056000in}}%
\pgfpathlineto{\pgfqpoint{1.872000in}{3.288000in}}%
\pgfusepath{fill}%
\end{pgfscope}%
\begin{pgfscope}%
\pgfpathrectangle{\pgfqpoint{1.200000in}{1.200000in}}{\pgfqpoint{3.360000in}{2.320000in}}%
\pgfusepath{clip}%
\pgfsetbuttcap%
\pgfsetroundjoin%
\definecolor{currentfill}{rgb}{0.336898,0.117591,0.308894}%
\pgfsetfillcolor{currentfill}%
\pgfsetlinewidth{0.000000pt}%
\definecolor{currentstroke}{rgb}{1.000000,1.000000,1.000000}%
\pgfsetstrokecolor{currentstroke}%
\pgfsetdash{}{0pt}%
\pgfpathmoveto{\pgfqpoint{2.208000in}{3.288000in}}%
\pgfpathlineto{\pgfqpoint{2.544000in}{3.288000in}}%
\pgfpathlineto{\pgfqpoint{2.544000in}{3.056000in}}%
\pgfpathlineto{\pgfqpoint{2.208000in}{3.056000in}}%
\pgfpathlineto{\pgfqpoint{2.208000in}{3.288000in}}%
\pgfusepath{fill}%
\end{pgfscope}%
\begin{pgfscope}%
\pgfpathrectangle{\pgfqpoint{1.200000in}{1.200000in}}{\pgfqpoint{3.360000in}{2.320000in}}%
\pgfusepath{clip}%
\pgfsetbuttcap%
\pgfsetroundjoin%
\definecolor{currentfill}{rgb}{0.458090,0.121430,0.345400}%
\pgfsetfillcolor{currentfill}%
\pgfsetlinewidth{0.000000pt}%
\definecolor{currentstroke}{rgb}{1.000000,1.000000,1.000000}%
\pgfsetstrokecolor{currentstroke}%
\pgfsetdash{}{0pt}%
\pgfpathmoveto{\pgfqpoint{2.544000in}{3.288000in}}%
\pgfpathlineto{\pgfqpoint{2.880000in}{3.288000in}}%
\pgfpathlineto{\pgfqpoint{2.880000in}{3.056000in}}%
\pgfpathlineto{\pgfqpoint{2.544000in}{3.056000in}}%
\pgfpathlineto{\pgfqpoint{2.544000in}{3.288000in}}%
\pgfusepath{fill}%
\end{pgfscope}%
\begin{pgfscope}%
\pgfpathrectangle{\pgfqpoint{1.200000in}{1.200000in}}{\pgfqpoint{3.360000in}{2.320000in}}%
\pgfusepath{clip}%
\pgfsetbuttcap%
\pgfsetroundjoin%
\definecolor{currentfill}{rgb}{0.380929,0.120615,0.325065}%
\pgfsetfillcolor{currentfill}%
\pgfsetlinewidth{0.000000pt}%
\definecolor{currentstroke}{rgb}{1.000000,1.000000,1.000000}%
\pgfsetstrokecolor{currentstroke}%
\pgfsetdash{}{0pt}%
\pgfpathmoveto{\pgfqpoint{2.880000in}{3.288000in}}%
\pgfpathlineto{\pgfqpoint{3.216000in}{3.288000in}}%
\pgfpathlineto{\pgfqpoint{3.216000in}{3.056000in}}%
\pgfpathlineto{\pgfqpoint{2.880000in}{3.056000in}}%
\pgfpathlineto{\pgfqpoint{2.880000in}{3.288000in}}%
\pgfusepath{fill}%
\end{pgfscope}%
\begin{pgfscope}%
\pgfpathrectangle{\pgfqpoint{1.200000in}{1.200000in}}{\pgfqpoint{3.360000in}{2.320000in}}%
\pgfusepath{clip}%
\pgfsetbuttcap%
\pgfsetroundjoin%
\definecolor{currentfill}{rgb}{0.361966,0.119541,0.318520}%
\pgfsetfillcolor{currentfill}%
\pgfsetlinewidth{0.000000pt}%
\definecolor{currentstroke}{rgb}{1.000000,1.000000,1.000000}%
\pgfsetstrokecolor{currentstroke}%
\pgfsetdash{}{0pt}%
\pgfpathmoveto{\pgfqpoint{3.216000in}{3.288000in}}%
\pgfpathlineto{\pgfqpoint{3.552000in}{3.288000in}}%
\pgfpathlineto{\pgfqpoint{3.552000in}{3.056000in}}%
\pgfpathlineto{\pgfqpoint{3.216000in}{3.056000in}}%
\pgfpathlineto{\pgfqpoint{3.216000in}{3.288000in}}%
\pgfusepath{fill}%
\end{pgfscope}%
\begin{pgfscope}%
\pgfpathrectangle{\pgfqpoint{1.200000in}{1.200000in}}{\pgfqpoint{3.360000in}{2.320000in}}%
\pgfusepath{clip}%
\pgfsetbuttcap%
\pgfsetroundjoin%
\definecolor{currentfill}{rgb}{0.361966,0.119541,0.318520}%
\pgfsetfillcolor{currentfill}%
\pgfsetlinewidth{0.000000pt}%
\definecolor{currentstroke}{rgb}{1.000000,1.000000,1.000000}%
\pgfsetstrokecolor{currentstroke}%
\pgfsetdash{}{0pt}%
\pgfpathmoveto{\pgfqpoint{3.552000in}{3.288000in}}%
\pgfpathlineto{\pgfqpoint{3.888000in}{3.288000in}}%
\pgfpathlineto{\pgfqpoint{3.888000in}{3.056000in}}%
\pgfpathlineto{\pgfqpoint{3.552000in}{3.056000in}}%
\pgfpathlineto{\pgfqpoint{3.552000in}{3.288000in}}%
\pgfusepath{fill}%
\end{pgfscope}%
\begin{pgfscope}%
\pgfpathrectangle{\pgfqpoint{1.200000in}{1.200000in}}{\pgfqpoint{3.360000in}{2.320000in}}%
\pgfusepath{clip}%
\pgfsetbuttcap%
\pgfsetroundjoin%
\definecolor{currentfill}{rgb}{0.336898,0.117591,0.308894}%
\pgfsetfillcolor{currentfill}%
\pgfsetlinewidth{0.000000pt}%
\definecolor{currentstroke}{rgb}{1.000000,1.000000,1.000000}%
\pgfsetstrokecolor{currentstroke}%
\pgfsetdash{}{0pt}%
\pgfpathmoveto{\pgfqpoint{3.888000in}{3.288000in}}%
\pgfpathlineto{\pgfqpoint{4.224000in}{3.288000in}}%
\pgfpathlineto{\pgfqpoint{4.224000in}{3.056000in}}%
\pgfpathlineto{\pgfqpoint{3.888000in}{3.056000in}}%
\pgfpathlineto{\pgfqpoint{3.888000in}{3.288000in}}%
\pgfusepath{fill}%
\end{pgfscope}%
\begin{pgfscope}%
\pgfpathrectangle{\pgfqpoint{1.200000in}{1.200000in}}{\pgfqpoint{3.360000in}{2.320000in}}%
\pgfusepath{clip}%
\pgfsetbuttcap%
\pgfsetroundjoin%
\definecolor{currentfill}{rgb}{0.380929,0.120615,0.325065}%
\pgfsetfillcolor{currentfill}%
\pgfsetlinewidth{0.000000pt}%
\definecolor{currentstroke}{rgb}{1.000000,1.000000,1.000000}%
\pgfsetstrokecolor{currentstroke}%
\pgfsetdash{}{0pt}%
\pgfpathmoveto{\pgfqpoint{4.224000in}{3.288000in}}%
\pgfpathlineto{\pgfqpoint{4.560000in}{3.288000in}}%
\pgfpathlineto{\pgfqpoint{4.560000in}{3.056000in}}%
\pgfpathlineto{\pgfqpoint{4.224000in}{3.056000in}}%
\pgfpathlineto{\pgfqpoint{4.224000in}{3.288000in}}%
\pgfusepath{fill}%
\end{pgfscope}%
\begin{pgfscope}%
\pgfpathrectangle{\pgfqpoint{1.200000in}{1.200000in}}{\pgfqpoint{3.360000in}{2.320000in}}%
\pgfusepath{clip}%
\pgfsetbuttcap%
\pgfsetroundjoin%
\definecolor{currentfill}{rgb}{0.336898,0.117591,0.308894}%
\pgfsetfillcolor{currentfill}%
\pgfsetlinewidth{0.000000pt}%
\definecolor{currentstroke}{rgb}{1.000000,1.000000,1.000000}%
\pgfsetstrokecolor{currentstroke}%
\pgfsetdash{}{0pt}%
\pgfpathmoveto{\pgfqpoint{1.200000in}{3.056000in}}%
\pgfpathlineto{\pgfqpoint{1.536000in}{3.056000in}}%
\pgfpathlineto{\pgfqpoint{1.536000in}{2.824000in}}%
\pgfpathlineto{\pgfqpoint{1.200000in}{2.824000in}}%
\pgfpathlineto{\pgfqpoint{1.200000in}{3.056000in}}%
\pgfusepath{fill}%
\end{pgfscope}%
\begin{pgfscope}%
\pgfpathrectangle{\pgfqpoint{1.200000in}{1.200000in}}{\pgfqpoint{3.360000in}{2.320000in}}%
\pgfusepath{clip}%
\pgfsetbuttcap%
\pgfsetroundjoin%
\definecolor{currentfill}{rgb}{0.318263,0.115743,0.301001}%
\pgfsetfillcolor{currentfill}%
\pgfsetlinewidth{0.000000pt}%
\definecolor{currentstroke}{rgb}{1.000000,1.000000,1.000000}%
\pgfsetstrokecolor{currentstroke}%
\pgfsetdash{}{0pt}%
\pgfpathmoveto{\pgfqpoint{1.536000in}{3.056000in}}%
\pgfpathlineto{\pgfqpoint{1.872000in}{3.056000in}}%
\pgfpathlineto{\pgfqpoint{1.872000in}{2.824000in}}%
\pgfpathlineto{\pgfqpoint{1.536000in}{2.824000in}}%
\pgfpathlineto{\pgfqpoint{1.536000in}{3.056000in}}%
\pgfusepath{fill}%
\end{pgfscope}%
\begin{pgfscope}%
\pgfpathrectangle{\pgfqpoint{1.200000in}{1.200000in}}{\pgfqpoint{3.360000in}{2.320000in}}%
\pgfusepath{clip}%
\pgfsetbuttcap%
\pgfsetroundjoin%
\definecolor{currentfill}{rgb}{0.981377,0.920617,0.865369}%
\pgfsetfillcolor{currentfill}%
\pgfsetlinewidth{0.000000pt}%
\definecolor{currentstroke}{rgb}{1.000000,1.000000,1.000000}%
\pgfsetstrokecolor{currentstroke}%
\pgfsetdash{}{0pt}%
\pgfpathmoveto{\pgfqpoint{1.872000in}{3.056000in}}%
\pgfpathlineto{\pgfqpoint{2.208000in}{3.056000in}}%
\pgfpathlineto{\pgfqpoint{2.208000in}{2.824000in}}%
\pgfpathlineto{\pgfqpoint{1.872000in}{2.824000in}}%
\pgfpathlineto{\pgfqpoint{1.872000in}{3.056000in}}%
\pgfusepath{fill}%
\end{pgfscope}%
\begin{pgfscope}%
\pgfpathrectangle{\pgfqpoint{1.200000in}{1.200000in}}{\pgfqpoint{3.360000in}{2.320000in}}%
\pgfusepath{clip}%
\pgfsetbuttcap%
\pgfsetroundjoin%
\definecolor{currentfill}{rgb}{0.380929,0.120615,0.325065}%
\pgfsetfillcolor{currentfill}%
\pgfsetlinewidth{0.000000pt}%
\definecolor{currentstroke}{rgb}{1.000000,1.000000,1.000000}%
\pgfsetstrokecolor{currentstroke}%
\pgfsetdash{}{0pt}%
\pgfpathmoveto{\pgfqpoint{2.208000in}{3.056000in}}%
\pgfpathlineto{\pgfqpoint{2.544000in}{3.056000in}}%
\pgfpathlineto{\pgfqpoint{2.544000in}{2.824000in}}%
\pgfpathlineto{\pgfqpoint{2.208000in}{2.824000in}}%
\pgfpathlineto{\pgfqpoint{2.208000in}{3.056000in}}%
\pgfusepath{fill}%
\end{pgfscope}%
\begin{pgfscope}%
\pgfpathrectangle{\pgfqpoint{1.200000in}{1.200000in}}{\pgfqpoint{3.360000in}{2.320000in}}%
\pgfusepath{clip}%
\pgfsetbuttcap%
\pgfsetroundjoin%
\definecolor{currentfill}{rgb}{0.336898,0.117591,0.308894}%
\pgfsetfillcolor{currentfill}%
\pgfsetlinewidth{0.000000pt}%
\definecolor{currentstroke}{rgb}{1.000000,1.000000,1.000000}%
\pgfsetstrokecolor{currentstroke}%
\pgfsetdash{}{0pt}%
\pgfpathmoveto{\pgfqpoint{2.544000in}{3.056000in}}%
\pgfpathlineto{\pgfqpoint{2.880000in}{3.056000in}}%
\pgfpathlineto{\pgfqpoint{2.880000in}{2.824000in}}%
\pgfpathlineto{\pgfqpoint{2.544000in}{2.824000in}}%
\pgfpathlineto{\pgfqpoint{2.544000in}{3.056000in}}%
\pgfusepath{fill}%
\end{pgfscope}%
\begin{pgfscope}%
\pgfpathrectangle{\pgfqpoint{1.200000in}{1.200000in}}{\pgfqpoint{3.360000in}{2.320000in}}%
\pgfusepath{clip}%
\pgfsetbuttcap%
\pgfsetroundjoin%
\definecolor{currentfill}{rgb}{0.368272,0.119937,0.320767}%
\pgfsetfillcolor{currentfill}%
\pgfsetlinewidth{0.000000pt}%
\definecolor{currentstroke}{rgb}{1.000000,1.000000,1.000000}%
\pgfsetstrokecolor{currentstroke}%
\pgfsetdash{}{0pt}%
\pgfpathmoveto{\pgfqpoint{2.880000in}{3.056000in}}%
\pgfpathlineto{\pgfqpoint{3.216000in}{3.056000in}}%
\pgfpathlineto{\pgfqpoint{3.216000in}{2.824000in}}%
\pgfpathlineto{\pgfqpoint{2.880000in}{2.824000in}}%
\pgfpathlineto{\pgfqpoint{2.880000in}{3.056000in}}%
\pgfusepath{fill}%
\end{pgfscope}%
\begin{pgfscope}%
\pgfpathrectangle{\pgfqpoint{1.200000in}{1.200000in}}{\pgfqpoint{3.360000in}{2.320000in}}%
\pgfusepath{clip}%
\pgfsetbuttcap%
\pgfsetroundjoin%
\definecolor{currentfill}{rgb}{0.368272,0.119937,0.320767}%
\pgfsetfillcolor{currentfill}%
\pgfsetlinewidth{0.000000pt}%
\definecolor{currentstroke}{rgb}{1.000000,1.000000,1.000000}%
\pgfsetstrokecolor{currentstroke}%
\pgfsetdash{}{0pt}%
\pgfpathmoveto{\pgfqpoint{3.216000in}{3.056000in}}%
\pgfpathlineto{\pgfqpoint{3.552000in}{3.056000in}}%
\pgfpathlineto{\pgfqpoint{3.552000in}{2.824000in}}%
\pgfpathlineto{\pgfqpoint{3.216000in}{2.824000in}}%
\pgfpathlineto{\pgfqpoint{3.216000in}{3.056000in}}%
\pgfusepath{fill}%
\end{pgfscope}%
\begin{pgfscope}%
\pgfpathrectangle{\pgfqpoint{1.200000in}{1.200000in}}{\pgfqpoint{3.360000in}{2.320000in}}%
\pgfusepath{clip}%
\pgfsetbuttcap%
\pgfsetroundjoin%
\definecolor{currentfill}{rgb}{0.355676,0.119109,0.316210}%
\pgfsetfillcolor{currentfill}%
\pgfsetlinewidth{0.000000pt}%
\definecolor{currentstroke}{rgb}{1.000000,1.000000,1.000000}%
\pgfsetstrokecolor{currentstroke}%
\pgfsetdash{}{0pt}%
\pgfpathmoveto{\pgfqpoint{3.552000in}{3.056000in}}%
\pgfpathlineto{\pgfqpoint{3.888000in}{3.056000in}}%
\pgfpathlineto{\pgfqpoint{3.888000in}{2.824000in}}%
\pgfpathlineto{\pgfqpoint{3.552000in}{2.824000in}}%
\pgfpathlineto{\pgfqpoint{3.552000in}{3.056000in}}%
\pgfusepath{fill}%
\end{pgfscope}%
\begin{pgfscope}%
\pgfpathrectangle{\pgfqpoint{1.200000in}{1.200000in}}{\pgfqpoint{3.360000in}{2.320000in}}%
\pgfusepath{clip}%
\pgfsetbuttcap%
\pgfsetroundjoin%
\definecolor{currentfill}{rgb}{0.361966,0.119541,0.318520}%
\pgfsetfillcolor{currentfill}%
\pgfsetlinewidth{0.000000pt}%
\definecolor{currentstroke}{rgb}{1.000000,1.000000,1.000000}%
\pgfsetstrokecolor{currentstroke}%
\pgfsetdash{}{0pt}%
\pgfpathmoveto{\pgfqpoint{3.888000in}{3.056000in}}%
\pgfpathlineto{\pgfqpoint{4.224000in}{3.056000in}}%
\pgfpathlineto{\pgfqpoint{4.224000in}{2.824000in}}%
\pgfpathlineto{\pgfqpoint{3.888000in}{2.824000in}}%
\pgfpathlineto{\pgfqpoint{3.888000in}{3.056000in}}%
\pgfusepath{fill}%
\end{pgfscope}%
\begin{pgfscope}%
\pgfpathrectangle{\pgfqpoint{1.200000in}{1.200000in}}{\pgfqpoint{3.360000in}{2.320000in}}%
\pgfusepath{clip}%
\pgfsetbuttcap%
\pgfsetroundjoin%
\definecolor{currentfill}{rgb}{0.349401,0.118640,0.313835}%
\pgfsetfillcolor{currentfill}%
\pgfsetlinewidth{0.000000pt}%
\definecolor{currentstroke}{rgb}{1.000000,1.000000,1.000000}%
\pgfsetstrokecolor{currentstroke}%
\pgfsetdash{}{0pt}%
\pgfpathmoveto{\pgfqpoint{4.224000in}{3.056000in}}%
\pgfpathlineto{\pgfqpoint{4.560000in}{3.056000in}}%
\pgfpathlineto{\pgfqpoint{4.560000in}{2.824000in}}%
\pgfpathlineto{\pgfqpoint{4.224000in}{2.824000in}}%
\pgfpathlineto{\pgfqpoint{4.224000in}{3.056000in}}%
\pgfusepath{fill}%
\end{pgfscope}%
\begin{pgfscope}%
\pgfpathrectangle{\pgfqpoint{1.200000in}{1.200000in}}{\pgfqpoint{3.360000in}{2.320000in}}%
\pgfusepath{clip}%
\pgfsetbuttcap%
\pgfsetroundjoin%
\definecolor{currentfill}{rgb}{0.330670,0.117012,0.306326}%
\pgfsetfillcolor{currentfill}%
\pgfsetlinewidth{0.000000pt}%
\definecolor{currentstroke}{rgb}{1.000000,1.000000,1.000000}%
\pgfsetstrokecolor{currentstroke}%
\pgfsetdash{}{0pt}%
\pgfpathmoveto{\pgfqpoint{1.200000in}{2.824000in}}%
\pgfpathlineto{\pgfqpoint{1.536000in}{2.824000in}}%
\pgfpathlineto{\pgfqpoint{1.536000in}{2.592000in}}%
\pgfpathlineto{\pgfqpoint{1.200000in}{2.592000in}}%
\pgfpathlineto{\pgfqpoint{1.200000in}{2.824000in}}%
\pgfusepath{fill}%
\end{pgfscope}%
\begin{pgfscope}%
\pgfpathrectangle{\pgfqpoint{1.200000in}{1.200000in}}{\pgfqpoint{3.360000in}{2.320000in}}%
\pgfusepath{clip}%
\pgfsetbuttcap%
\pgfsetroundjoin%
\definecolor{currentfill}{rgb}{0.336898,0.117591,0.308894}%
\pgfsetfillcolor{currentfill}%
\pgfsetlinewidth{0.000000pt}%
\definecolor{currentstroke}{rgb}{1.000000,1.000000,1.000000}%
\pgfsetstrokecolor{currentstroke}%
\pgfsetdash{}{0pt}%
\pgfpathmoveto{\pgfqpoint{1.536000in}{2.824000in}}%
\pgfpathlineto{\pgfqpoint{1.872000in}{2.824000in}}%
\pgfpathlineto{\pgfqpoint{1.872000in}{2.592000in}}%
\pgfpathlineto{\pgfqpoint{1.536000in}{2.592000in}}%
\pgfpathlineto{\pgfqpoint{1.536000in}{2.824000in}}%
\pgfusepath{fill}%
\end{pgfscope}%
\begin{pgfscope}%
\pgfpathrectangle{\pgfqpoint{1.200000in}{1.200000in}}{\pgfqpoint{3.360000in}{2.320000in}}%
\pgfusepath{clip}%
\pgfsetbuttcap%
\pgfsetroundjoin%
\definecolor{currentfill}{rgb}{0.380929,0.120615,0.325065}%
\pgfsetfillcolor{currentfill}%
\pgfsetlinewidth{0.000000pt}%
\definecolor{currentstroke}{rgb}{1.000000,1.000000,1.000000}%
\pgfsetstrokecolor{currentstroke}%
\pgfsetdash{}{0pt}%
\pgfpathmoveto{\pgfqpoint{1.872000in}{2.824000in}}%
\pgfpathlineto{\pgfqpoint{2.208000in}{2.824000in}}%
\pgfpathlineto{\pgfqpoint{2.208000in}{2.592000in}}%
\pgfpathlineto{\pgfqpoint{1.872000in}{2.592000in}}%
\pgfpathlineto{\pgfqpoint{1.872000in}{2.824000in}}%
\pgfusepath{fill}%
\end{pgfscope}%
\begin{pgfscope}%
\pgfpathrectangle{\pgfqpoint{1.200000in}{1.200000in}}{\pgfqpoint{3.360000in}{2.320000in}}%
\pgfusepath{clip}%
\pgfsetbuttcap%
\pgfsetroundjoin%
\definecolor{currentfill}{rgb}{0.981377,0.920617,0.865369}%
\pgfsetfillcolor{currentfill}%
\pgfsetlinewidth{0.000000pt}%
\definecolor{currentstroke}{rgb}{1.000000,1.000000,1.000000}%
\pgfsetstrokecolor{currentstroke}%
\pgfsetdash{}{0pt}%
\pgfpathmoveto{\pgfqpoint{2.208000in}{2.824000in}}%
\pgfpathlineto{\pgfqpoint{2.544000in}{2.824000in}}%
\pgfpathlineto{\pgfqpoint{2.544000in}{2.592000in}}%
\pgfpathlineto{\pgfqpoint{2.208000in}{2.592000in}}%
\pgfpathlineto{\pgfqpoint{2.208000in}{2.824000in}}%
\pgfusepath{fill}%
\end{pgfscope}%
\begin{pgfscope}%
\pgfpathrectangle{\pgfqpoint{1.200000in}{1.200000in}}{\pgfqpoint{3.360000in}{2.320000in}}%
\pgfusepath{clip}%
\pgfsetbuttcap%
\pgfsetroundjoin%
\definecolor{currentfill}{rgb}{0.336898,0.117591,0.308894}%
\pgfsetfillcolor{currentfill}%
\pgfsetlinewidth{0.000000pt}%
\definecolor{currentstroke}{rgb}{1.000000,1.000000,1.000000}%
\pgfsetstrokecolor{currentstroke}%
\pgfsetdash{}{0pt}%
\pgfpathmoveto{\pgfqpoint{2.544000in}{2.824000in}}%
\pgfpathlineto{\pgfqpoint{2.880000in}{2.824000in}}%
\pgfpathlineto{\pgfqpoint{2.880000in}{2.592000in}}%
\pgfpathlineto{\pgfqpoint{2.544000in}{2.592000in}}%
\pgfpathlineto{\pgfqpoint{2.544000in}{2.824000in}}%
\pgfusepath{fill}%
\end{pgfscope}%
\begin{pgfscope}%
\pgfpathrectangle{\pgfqpoint{1.200000in}{1.200000in}}{\pgfqpoint{3.360000in}{2.320000in}}%
\pgfusepath{clip}%
\pgfsetbuttcap%
\pgfsetroundjoin%
\definecolor{currentfill}{rgb}{0.355676,0.119109,0.316210}%
\pgfsetfillcolor{currentfill}%
\pgfsetlinewidth{0.000000pt}%
\definecolor{currentstroke}{rgb}{1.000000,1.000000,1.000000}%
\pgfsetstrokecolor{currentstroke}%
\pgfsetdash{}{0pt}%
\pgfpathmoveto{\pgfqpoint{2.880000in}{2.824000in}}%
\pgfpathlineto{\pgfqpoint{3.216000in}{2.824000in}}%
\pgfpathlineto{\pgfqpoint{3.216000in}{2.592000in}}%
\pgfpathlineto{\pgfqpoint{2.880000in}{2.592000in}}%
\pgfpathlineto{\pgfqpoint{2.880000in}{2.824000in}}%
\pgfusepath{fill}%
\end{pgfscope}%
\begin{pgfscope}%
\pgfpathrectangle{\pgfqpoint{1.200000in}{1.200000in}}{\pgfqpoint{3.360000in}{2.320000in}}%
\pgfusepath{clip}%
\pgfsetbuttcap%
\pgfsetroundjoin%
\definecolor{currentfill}{rgb}{0.355676,0.119109,0.316210}%
\pgfsetfillcolor{currentfill}%
\pgfsetlinewidth{0.000000pt}%
\definecolor{currentstroke}{rgb}{1.000000,1.000000,1.000000}%
\pgfsetstrokecolor{currentstroke}%
\pgfsetdash{}{0pt}%
\pgfpathmoveto{\pgfqpoint{3.216000in}{2.824000in}}%
\pgfpathlineto{\pgfqpoint{3.552000in}{2.824000in}}%
\pgfpathlineto{\pgfqpoint{3.552000in}{2.592000in}}%
\pgfpathlineto{\pgfqpoint{3.216000in}{2.592000in}}%
\pgfpathlineto{\pgfqpoint{3.216000in}{2.824000in}}%
\pgfusepath{fill}%
\end{pgfscope}%
\begin{pgfscope}%
\pgfpathrectangle{\pgfqpoint{1.200000in}{1.200000in}}{\pgfqpoint{3.360000in}{2.320000in}}%
\pgfusepath{clip}%
\pgfsetbuttcap%
\pgfsetroundjoin%
\definecolor{currentfill}{rgb}{0.343142,0.118134,0.311397}%
\pgfsetfillcolor{currentfill}%
\pgfsetlinewidth{0.000000pt}%
\definecolor{currentstroke}{rgb}{1.000000,1.000000,1.000000}%
\pgfsetstrokecolor{currentstroke}%
\pgfsetdash{}{0pt}%
\pgfpathmoveto{\pgfqpoint{3.552000in}{2.824000in}}%
\pgfpathlineto{\pgfqpoint{3.888000in}{2.824000in}}%
\pgfpathlineto{\pgfqpoint{3.888000in}{2.592000in}}%
\pgfpathlineto{\pgfqpoint{3.552000in}{2.592000in}}%
\pgfpathlineto{\pgfqpoint{3.552000in}{2.824000in}}%
\pgfusepath{fill}%
\end{pgfscope}%
\begin{pgfscope}%
\pgfpathrectangle{\pgfqpoint{1.200000in}{1.200000in}}{\pgfqpoint{3.360000in}{2.320000in}}%
\pgfusepath{clip}%
\pgfsetbuttcap%
\pgfsetroundjoin%
\definecolor{currentfill}{rgb}{0.336898,0.117591,0.308894}%
\pgfsetfillcolor{currentfill}%
\pgfsetlinewidth{0.000000pt}%
\definecolor{currentstroke}{rgb}{1.000000,1.000000,1.000000}%
\pgfsetstrokecolor{currentstroke}%
\pgfsetdash{}{0pt}%
\pgfpathmoveto{\pgfqpoint{3.888000in}{2.824000in}}%
\pgfpathlineto{\pgfqpoint{4.224000in}{2.824000in}}%
\pgfpathlineto{\pgfqpoint{4.224000in}{2.592000in}}%
\pgfpathlineto{\pgfqpoint{3.888000in}{2.592000in}}%
\pgfpathlineto{\pgfqpoint{3.888000in}{2.824000in}}%
\pgfusepath{fill}%
\end{pgfscope}%
\begin{pgfscope}%
\pgfpathrectangle{\pgfqpoint{1.200000in}{1.200000in}}{\pgfqpoint{3.360000in}{2.320000in}}%
\pgfusepath{clip}%
\pgfsetbuttcap%
\pgfsetroundjoin%
\definecolor{currentfill}{rgb}{0.343142,0.118134,0.311397}%
\pgfsetfillcolor{currentfill}%
\pgfsetlinewidth{0.000000pt}%
\definecolor{currentstroke}{rgb}{1.000000,1.000000,1.000000}%
\pgfsetstrokecolor{currentstroke}%
\pgfsetdash{}{0pt}%
\pgfpathmoveto{\pgfqpoint{4.224000in}{2.824000in}}%
\pgfpathlineto{\pgfqpoint{4.560000in}{2.824000in}}%
\pgfpathlineto{\pgfqpoint{4.560000in}{2.592000in}}%
\pgfpathlineto{\pgfqpoint{4.224000in}{2.592000in}}%
\pgfpathlineto{\pgfqpoint{4.224000in}{2.824000in}}%
\pgfusepath{fill}%
\end{pgfscope}%
\begin{pgfscope}%
\pgfpathrectangle{\pgfqpoint{1.200000in}{1.200000in}}{\pgfqpoint{3.360000in}{2.320000in}}%
\pgfusepath{clip}%
\pgfsetbuttcap%
\pgfsetroundjoin%
\definecolor{currentfill}{rgb}{0.387280,0.120898,0.327118}%
\pgfsetfillcolor{currentfill}%
\pgfsetlinewidth{0.000000pt}%
\definecolor{currentstroke}{rgb}{1.000000,1.000000,1.000000}%
\pgfsetstrokecolor{currentstroke}%
\pgfsetdash{}{0pt}%
\pgfpathmoveto{\pgfqpoint{1.200000in}{2.592000in}}%
\pgfpathlineto{\pgfqpoint{1.536000in}{2.592000in}}%
\pgfpathlineto{\pgfqpoint{1.536000in}{2.360000in}}%
\pgfpathlineto{\pgfqpoint{1.200000in}{2.360000in}}%
\pgfpathlineto{\pgfqpoint{1.200000in}{2.592000in}}%
\pgfusepath{fill}%
\end{pgfscope}%
\begin{pgfscope}%
\pgfpathrectangle{\pgfqpoint{1.200000in}{1.200000in}}{\pgfqpoint{3.360000in}{2.320000in}}%
\pgfusepath{clip}%
\pgfsetbuttcap%
\pgfsetroundjoin%
\definecolor{currentfill}{rgb}{0.458090,0.121430,0.345400}%
\pgfsetfillcolor{currentfill}%
\pgfsetlinewidth{0.000000pt}%
\definecolor{currentstroke}{rgb}{1.000000,1.000000,1.000000}%
\pgfsetstrokecolor{currentstroke}%
\pgfsetdash{}{0pt}%
\pgfpathmoveto{\pgfqpoint{1.536000in}{2.592000in}}%
\pgfpathlineto{\pgfqpoint{1.872000in}{2.592000in}}%
\pgfpathlineto{\pgfqpoint{1.872000in}{2.360000in}}%
\pgfpathlineto{\pgfqpoint{1.536000in}{2.360000in}}%
\pgfpathlineto{\pgfqpoint{1.536000in}{2.592000in}}%
\pgfusepath{fill}%
\end{pgfscope}%
\begin{pgfscope}%
\pgfpathrectangle{\pgfqpoint{1.200000in}{1.200000in}}{\pgfqpoint{3.360000in}{2.320000in}}%
\pgfusepath{clip}%
\pgfsetbuttcap%
\pgfsetroundjoin%
\definecolor{currentfill}{rgb}{0.336898,0.117591,0.308894}%
\pgfsetfillcolor{currentfill}%
\pgfsetlinewidth{0.000000pt}%
\definecolor{currentstroke}{rgb}{1.000000,1.000000,1.000000}%
\pgfsetstrokecolor{currentstroke}%
\pgfsetdash{}{0pt}%
\pgfpathmoveto{\pgfqpoint{1.872000in}{2.592000in}}%
\pgfpathlineto{\pgfqpoint{2.208000in}{2.592000in}}%
\pgfpathlineto{\pgfqpoint{2.208000in}{2.360000in}}%
\pgfpathlineto{\pgfqpoint{1.872000in}{2.360000in}}%
\pgfpathlineto{\pgfqpoint{1.872000in}{2.592000in}}%
\pgfusepath{fill}%
\end{pgfscope}%
\begin{pgfscope}%
\pgfpathrectangle{\pgfqpoint{1.200000in}{1.200000in}}{\pgfqpoint{3.360000in}{2.320000in}}%
\pgfusepath{clip}%
\pgfsetbuttcap%
\pgfsetroundjoin%
\definecolor{currentfill}{rgb}{0.336898,0.117591,0.308894}%
\pgfsetfillcolor{currentfill}%
\pgfsetlinewidth{0.000000pt}%
\definecolor{currentstroke}{rgb}{1.000000,1.000000,1.000000}%
\pgfsetstrokecolor{currentstroke}%
\pgfsetdash{}{0pt}%
\pgfpathmoveto{\pgfqpoint{2.208000in}{2.592000in}}%
\pgfpathlineto{\pgfqpoint{2.544000in}{2.592000in}}%
\pgfpathlineto{\pgfqpoint{2.544000in}{2.360000in}}%
\pgfpathlineto{\pgfqpoint{2.208000in}{2.360000in}}%
\pgfpathlineto{\pgfqpoint{2.208000in}{2.592000in}}%
\pgfusepath{fill}%
\end{pgfscope}%
\begin{pgfscope}%
\pgfpathrectangle{\pgfqpoint{1.200000in}{1.200000in}}{\pgfqpoint{3.360000in}{2.320000in}}%
\pgfusepath{clip}%
\pgfsetbuttcap%
\pgfsetroundjoin%
\definecolor{currentfill}{rgb}{0.981377,0.920617,0.865369}%
\pgfsetfillcolor{currentfill}%
\pgfsetlinewidth{0.000000pt}%
\definecolor{currentstroke}{rgb}{1.000000,1.000000,1.000000}%
\pgfsetstrokecolor{currentstroke}%
\pgfsetdash{}{0pt}%
\pgfpathmoveto{\pgfqpoint{2.544000in}{2.592000in}}%
\pgfpathlineto{\pgfqpoint{2.880000in}{2.592000in}}%
\pgfpathlineto{\pgfqpoint{2.880000in}{2.360000in}}%
\pgfpathlineto{\pgfqpoint{2.544000in}{2.360000in}}%
\pgfpathlineto{\pgfqpoint{2.544000in}{2.592000in}}%
\pgfusepath{fill}%
\end{pgfscope}%
\begin{pgfscope}%
\pgfpathrectangle{\pgfqpoint{1.200000in}{1.200000in}}{\pgfqpoint{3.360000in}{2.320000in}}%
\pgfusepath{clip}%
\pgfsetbuttcap%
\pgfsetroundjoin%
\definecolor{currentfill}{rgb}{0.318263,0.115743,0.301001}%
\pgfsetfillcolor{currentfill}%
\pgfsetlinewidth{0.000000pt}%
\definecolor{currentstroke}{rgb}{1.000000,1.000000,1.000000}%
\pgfsetstrokecolor{currentstroke}%
\pgfsetdash{}{0pt}%
\pgfpathmoveto{\pgfqpoint{2.880000in}{2.592000in}}%
\pgfpathlineto{\pgfqpoint{3.216000in}{2.592000in}}%
\pgfpathlineto{\pgfqpoint{3.216000in}{2.360000in}}%
\pgfpathlineto{\pgfqpoint{2.880000in}{2.360000in}}%
\pgfpathlineto{\pgfqpoint{2.880000in}{2.592000in}}%
\pgfusepath{fill}%
\end{pgfscope}%
\begin{pgfscope}%
\pgfpathrectangle{\pgfqpoint{1.200000in}{1.200000in}}{\pgfqpoint{3.360000in}{2.320000in}}%
\pgfusepath{clip}%
\pgfsetbuttcap%
\pgfsetroundjoin%
\definecolor{currentfill}{rgb}{0.312084,0.115053,0.298245}%
\pgfsetfillcolor{currentfill}%
\pgfsetlinewidth{0.000000pt}%
\definecolor{currentstroke}{rgb}{1.000000,1.000000,1.000000}%
\pgfsetstrokecolor{currentstroke}%
\pgfsetdash{}{0pt}%
\pgfpathmoveto{\pgfqpoint{3.216000in}{2.592000in}}%
\pgfpathlineto{\pgfqpoint{3.552000in}{2.592000in}}%
\pgfpathlineto{\pgfqpoint{3.552000in}{2.360000in}}%
\pgfpathlineto{\pgfqpoint{3.216000in}{2.360000in}}%
\pgfpathlineto{\pgfqpoint{3.216000in}{2.592000in}}%
\pgfusepath{fill}%
\end{pgfscope}%
\begin{pgfscope}%
\pgfpathrectangle{\pgfqpoint{1.200000in}{1.200000in}}{\pgfqpoint{3.360000in}{2.320000in}}%
\pgfusepath{clip}%
\pgfsetbuttcap%
\pgfsetroundjoin%
\definecolor{currentfill}{rgb}{0.380929,0.120615,0.325065}%
\pgfsetfillcolor{currentfill}%
\pgfsetlinewidth{0.000000pt}%
\definecolor{currentstroke}{rgb}{1.000000,1.000000,1.000000}%
\pgfsetstrokecolor{currentstroke}%
\pgfsetdash{}{0pt}%
\pgfpathmoveto{\pgfqpoint{3.552000in}{2.592000in}}%
\pgfpathlineto{\pgfqpoint{3.888000in}{2.592000in}}%
\pgfpathlineto{\pgfqpoint{3.888000in}{2.360000in}}%
\pgfpathlineto{\pgfqpoint{3.552000in}{2.360000in}}%
\pgfpathlineto{\pgfqpoint{3.552000in}{2.592000in}}%
\pgfusepath{fill}%
\end{pgfscope}%
\begin{pgfscope}%
\pgfpathrectangle{\pgfqpoint{1.200000in}{1.200000in}}{\pgfqpoint{3.360000in}{2.320000in}}%
\pgfusepath{clip}%
\pgfsetbuttcap%
\pgfsetroundjoin%
\definecolor{currentfill}{rgb}{0.343142,0.118134,0.311397}%
\pgfsetfillcolor{currentfill}%
\pgfsetlinewidth{0.000000pt}%
\definecolor{currentstroke}{rgb}{1.000000,1.000000,1.000000}%
\pgfsetstrokecolor{currentstroke}%
\pgfsetdash{}{0pt}%
\pgfpathmoveto{\pgfqpoint{3.888000in}{2.592000in}}%
\pgfpathlineto{\pgfqpoint{4.224000in}{2.592000in}}%
\pgfpathlineto{\pgfqpoint{4.224000in}{2.360000in}}%
\pgfpathlineto{\pgfqpoint{3.888000in}{2.360000in}}%
\pgfpathlineto{\pgfqpoint{3.888000in}{2.592000in}}%
\pgfusepath{fill}%
\end{pgfscope}%
\begin{pgfscope}%
\pgfpathrectangle{\pgfqpoint{1.200000in}{1.200000in}}{\pgfqpoint{3.360000in}{2.320000in}}%
\pgfusepath{clip}%
\pgfsetbuttcap%
\pgfsetroundjoin%
\definecolor{currentfill}{rgb}{0.343142,0.118134,0.311397}%
\pgfsetfillcolor{currentfill}%
\pgfsetlinewidth{0.000000pt}%
\definecolor{currentstroke}{rgb}{1.000000,1.000000,1.000000}%
\pgfsetstrokecolor{currentstroke}%
\pgfsetdash{}{0pt}%
\pgfpathmoveto{\pgfqpoint{4.224000in}{2.592000in}}%
\pgfpathlineto{\pgfqpoint{4.560000in}{2.592000in}}%
\pgfpathlineto{\pgfqpoint{4.560000in}{2.360000in}}%
\pgfpathlineto{\pgfqpoint{4.224000in}{2.360000in}}%
\pgfpathlineto{\pgfqpoint{4.224000in}{2.592000in}}%
\pgfusepath{fill}%
\end{pgfscope}%
\begin{pgfscope}%
\pgfpathrectangle{\pgfqpoint{1.200000in}{1.200000in}}{\pgfqpoint{3.360000in}{2.320000in}}%
\pgfusepath{clip}%
\pgfsetbuttcap%
\pgfsetroundjoin%
\definecolor{currentfill}{rgb}{0.368272,0.119937,0.320767}%
\pgfsetfillcolor{currentfill}%
\pgfsetlinewidth{0.000000pt}%
\definecolor{currentstroke}{rgb}{1.000000,1.000000,1.000000}%
\pgfsetstrokecolor{currentstroke}%
\pgfsetdash{}{0pt}%
\pgfpathmoveto{\pgfqpoint{1.200000in}{2.360000in}}%
\pgfpathlineto{\pgfqpoint{1.536000in}{2.360000in}}%
\pgfpathlineto{\pgfqpoint{1.536000in}{2.128000in}}%
\pgfpathlineto{\pgfqpoint{1.200000in}{2.128000in}}%
\pgfpathlineto{\pgfqpoint{1.200000in}{2.360000in}}%
\pgfusepath{fill}%
\end{pgfscope}%
\begin{pgfscope}%
\pgfpathrectangle{\pgfqpoint{1.200000in}{1.200000in}}{\pgfqpoint{3.360000in}{2.320000in}}%
\pgfusepath{clip}%
\pgfsetbuttcap%
\pgfsetroundjoin%
\definecolor{currentfill}{rgb}{0.380929,0.120615,0.325065}%
\pgfsetfillcolor{currentfill}%
\pgfsetlinewidth{0.000000pt}%
\definecolor{currentstroke}{rgb}{1.000000,1.000000,1.000000}%
\pgfsetstrokecolor{currentstroke}%
\pgfsetdash{}{0pt}%
\pgfpathmoveto{\pgfqpoint{1.536000in}{2.360000in}}%
\pgfpathlineto{\pgfqpoint{1.872000in}{2.360000in}}%
\pgfpathlineto{\pgfqpoint{1.872000in}{2.128000in}}%
\pgfpathlineto{\pgfqpoint{1.536000in}{2.128000in}}%
\pgfpathlineto{\pgfqpoint{1.536000in}{2.360000in}}%
\pgfusepath{fill}%
\end{pgfscope}%
\begin{pgfscope}%
\pgfpathrectangle{\pgfqpoint{1.200000in}{1.200000in}}{\pgfqpoint{3.360000in}{2.320000in}}%
\pgfusepath{clip}%
\pgfsetbuttcap%
\pgfsetroundjoin%
\definecolor{currentfill}{rgb}{0.368272,0.119937,0.320767}%
\pgfsetfillcolor{currentfill}%
\pgfsetlinewidth{0.000000pt}%
\definecolor{currentstroke}{rgb}{1.000000,1.000000,1.000000}%
\pgfsetstrokecolor{currentstroke}%
\pgfsetdash{}{0pt}%
\pgfpathmoveto{\pgfqpoint{1.872000in}{2.360000in}}%
\pgfpathlineto{\pgfqpoint{2.208000in}{2.360000in}}%
\pgfpathlineto{\pgfqpoint{2.208000in}{2.128000in}}%
\pgfpathlineto{\pgfqpoint{1.872000in}{2.128000in}}%
\pgfpathlineto{\pgfqpoint{1.872000in}{2.360000in}}%
\pgfusepath{fill}%
\end{pgfscope}%
\begin{pgfscope}%
\pgfpathrectangle{\pgfqpoint{1.200000in}{1.200000in}}{\pgfqpoint{3.360000in}{2.320000in}}%
\pgfusepath{clip}%
\pgfsetbuttcap%
\pgfsetroundjoin%
\definecolor{currentfill}{rgb}{0.355676,0.119109,0.316210}%
\pgfsetfillcolor{currentfill}%
\pgfsetlinewidth{0.000000pt}%
\definecolor{currentstroke}{rgb}{1.000000,1.000000,1.000000}%
\pgfsetstrokecolor{currentstroke}%
\pgfsetdash{}{0pt}%
\pgfpathmoveto{\pgfqpoint{2.208000in}{2.360000in}}%
\pgfpathlineto{\pgfqpoint{2.544000in}{2.360000in}}%
\pgfpathlineto{\pgfqpoint{2.544000in}{2.128000in}}%
\pgfpathlineto{\pgfqpoint{2.208000in}{2.128000in}}%
\pgfpathlineto{\pgfqpoint{2.208000in}{2.360000in}}%
\pgfusepath{fill}%
\end{pgfscope}%
\begin{pgfscope}%
\pgfpathrectangle{\pgfqpoint{1.200000in}{1.200000in}}{\pgfqpoint{3.360000in}{2.320000in}}%
\pgfusepath{clip}%
\pgfsetbuttcap%
\pgfsetroundjoin%
\definecolor{currentfill}{rgb}{0.318263,0.115743,0.301001}%
\pgfsetfillcolor{currentfill}%
\pgfsetlinewidth{0.000000pt}%
\definecolor{currentstroke}{rgb}{1.000000,1.000000,1.000000}%
\pgfsetstrokecolor{currentstroke}%
\pgfsetdash{}{0pt}%
\pgfpathmoveto{\pgfqpoint{2.544000in}{2.360000in}}%
\pgfpathlineto{\pgfqpoint{2.880000in}{2.360000in}}%
\pgfpathlineto{\pgfqpoint{2.880000in}{2.128000in}}%
\pgfpathlineto{\pgfqpoint{2.544000in}{2.128000in}}%
\pgfpathlineto{\pgfqpoint{2.544000in}{2.360000in}}%
\pgfusepath{fill}%
\end{pgfscope}%
\begin{pgfscope}%
\pgfpathrectangle{\pgfqpoint{1.200000in}{1.200000in}}{\pgfqpoint{3.360000in}{2.320000in}}%
\pgfusepath{clip}%
\pgfsetbuttcap%
\pgfsetroundjoin%
\definecolor{currentfill}{rgb}{0.981377,0.920617,0.865369}%
\pgfsetfillcolor{currentfill}%
\pgfsetlinewidth{0.000000pt}%
\definecolor{currentstroke}{rgb}{1.000000,1.000000,1.000000}%
\pgfsetstrokecolor{currentstroke}%
\pgfsetdash{}{0pt}%
\pgfpathmoveto{\pgfqpoint{2.880000in}{2.360000in}}%
\pgfpathlineto{\pgfqpoint{3.216000in}{2.360000in}}%
\pgfpathlineto{\pgfqpoint{3.216000in}{2.128000in}}%
\pgfpathlineto{\pgfqpoint{2.880000in}{2.128000in}}%
\pgfpathlineto{\pgfqpoint{2.880000in}{2.360000in}}%
\pgfusepath{fill}%
\end{pgfscope}%
\begin{pgfscope}%
\pgfpathrectangle{\pgfqpoint{1.200000in}{1.200000in}}{\pgfqpoint{3.360000in}{2.320000in}}%
\pgfusepath{clip}%
\pgfsetbuttcap%
\pgfsetroundjoin%
\definecolor{currentfill}{rgb}{0.324459,0.116396,0.303694}%
\pgfsetfillcolor{currentfill}%
\pgfsetlinewidth{0.000000pt}%
\definecolor{currentstroke}{rgb}{1.000000,1.000000,1.000000}%
\pgfsetstrokecolor{currentstroke}%
\pgfsetdash{}{0pt}%
\pgfpathmoveto{\pgfqpoint{3.216000in}{2.360000in}}%
\pgfpathlineto{\pgfqpoint{3.552000in}{2.360000in}}%
\pgfpathlineto{\pgfqpoint{3.552000in}{2.128000in}}%
\pgfpathlineto{\pgfqpoint{3.216000in}{2.128000in}}%
\pgfpathlineto{\pgfqpoint{3.216000in}{2.360000in}}%
\pgfusepath{fill}%
\end{pgfscope}%
\begin{pgfscope}%
\pgfpathrectangle{\pgfqpoint{1.200000in}{1.200000in}}{\pgfqpoint{3.360000in}{2.320000in}}%
\pgfusepath{clip}%
\pgfsetbuttcap%
\pgfsetroundjoin%
\definecolor{currentfill}{rgb}{0.355676,0.119109,0.316210}%
\pgfsetfillcolor{currentfill}%
\pgfsetlinewidth{0.000000pt}%
\definecolor{currentstroke}{rgb}{1.000000,1.000000,1.000000}%
\pgfsetstrokecolor{currentstroke}%
\pgfsetdash{}{0pt}%
\pgfpathmoveto{\pgfqpoint{3.552000in}{2.360000in}}%
\pgfpathlineto{\pgfqpoint{3.888000in}{2.360000in}}%
\pgfpathlineto{\pgfqpoint{3.888000in}{2.128000in}}%
\pgfpathlineto{\pgfqpoint{3.552000in}{2.128000in}}%
\pgfpathlineto{\pgfqpoint{3.552000in}{2.360000in}}%
\pgfusepath{fill}%
\end{pgfscope}%
\begin{pgfscope}%
\pgfpathrectangle{\pgfqpoint{1.200000in}{1.200000in}}{\pgfqpoint{3.360000in}{2.320000in}}%
\pgfusepath{clip}%
\pgfsetbuttcap%
\pgfsetroundjoin%
\definecolor{currentfill}{rgb}{0.343142,0.118134,0.311397}%
\pgfsetfillcolor{currentfill}%
\pgfsetlinewidth{0.000000pt}%
\definecolor{currentstroke}{rgb}{1.000000,1.000000,1.000000}%
\pgfsetstrokecolor{currentstroke}%
\pgfsetdash{}{0pt}%
\pgfpathmoveto{\pgfqpoint{3.888000in}{2.360000in}}%
\pgfpathlineto{\pgfqpoint{4.224000in}{2.360000in}}%
\pgfpathlineto{\pgfqpoint{4.224000in}{2.128000in}}%
\pgfpathlineto{\pgfqpoint{3.888000in}{2.128000in}}%
\pgfpathlineto{\pgfqpoint{3.888000in}{2.360000in}}%
\pgfusepath{fill}%
\end{pgfscope}%
\begin{pgfscope}%
\pgfpathrectangle{\pgfqpoint{1.200000in}{1.200000in}}{\pgfqpoint{3.360000in}{2.320000in}}%
\pgfusepath{clip}%
\pgfsetbuttcap%
\pgfsetroundjoin%
\definecolor{currentfill}{rgb}{0.349401,0.118640,0.313835}%
\pgfsetfillcolor{currentfill}%
\pgfsetlinewidth{0.000000pt}%
\definecolor{currentstroke}{rgb}{1.000000,1.000000,1.000000}%
\pgfsetstrokecolor{currentstroke}%
\pgfsetdash{}{0pt}%
\pgfpathmoveto{\pgfqpoint{4.224000in}{2.360000in}}%
\pgfpathlineto{\pgfqpoint{4.560000in}{2.360000in}}%
\pgfpathlineto{\pgfqpoint{4.560000in}{2.128000in}}%
\pgfpathlineto{\pgfqpoint{4.224000in}{2.128000in}}%
\pgfpathlineto{\pgfqpoint{4.224000in}{2.360000in}}%
\pgfusepath{fill}%
\end{pgfscope}%
\begin{pgfscope}%
\pgfpathrectangle{\pgfqpoint{1.200000in}{1.200000in}}{\pgfqpoint{3.360000in}{2.320000in}}%
\pgfusepath{clip}%
\pgfsetbuttcap%
\pgfsetroundjoin%
\definecolor{currentfill}{rgb}{0.010608,0.018082,0.100187}%
\pgfsetfillcolor{currentfill}%
\pgfsetlinewidth{0.000000pt}%
\definecolor{currentstroke}{rgb}{1.000000,1.000000,1.000000}%
\pgfsetstrokecolor{currentstroke}%
\pgfsetdash{}{0pt}%
\pgfpathmoveto{\pgfqpoint{1.200000in}{2.128000in}}%
\pgfpathlineto{\pgfqpoint{1.536000in}{2.128000in}}%
\pgfpathlineto{\pgfqpoint{1.536000in}{1.896000in}}%
\pgfpathlineto{\pgfqpoint{1.200000in}{1.896000in}}%
\pgfpathlineto{\pgfqpoint{1.200000in}{2.128000in}}%
\pgfusepath{fill}%
\end{pgfscope}%
\begin{pgfscope}%
\pgfpathrectangle{\pgfqpoint{1.200000in}{1.200000in}}{\pgfqpoint{3.360000in}{2.320000in}}%
\pgfusepath{clip}%
\pgfsetbuttcap%
\pgfsetroundjoin%
\definecolor{currentfill}{rgb}{0.361966,0.119541,0.318520}%
\pgfsetfillcolor{currentfill}%
\pgfsetlinewidth{0.000000pt}%
\definecolor{currentstroke}{rgb}{1.000000,1.000000,1.000000}%
\pgfsetstrokecolor{currentstroke}%
\pgfsetdash{}{0pt}%
\pgfpathmoveto{\pgfqpoint{1.536000in}{2.128000in}}%
\pgfpathlineto{\pgfqpoint{1.872000in}{2.128000in}}%
\pgfpathlineto{\pgfqpoint{1.872000in}{1.896000in}}%
\pgfpathlineto{\pgfqpoint{1.536000in}{1.896000in}}%
\pgfpathlineto{\pgfqpoint{1.536000in}{2.128000in}}%
\pgfusepath{fill}%
\end{pgfscope}%
\begin{pgfscope}%
\pgfpathrectangle{\pgfqpoint{1.200000in}{1.200000in}}{\pgfqpoint{3.360000in}{2.320000in}}%
\pgfusepath{clip}%
\pgfsetbuttcap%
\pgfsetroundjoin%
\definecolor{currentfill}{rgb}{0.368272,0.119937,0.320767}%
\pgfsetfillcolor{currentfill}%
\pgfsetlinewidth{0.000000pt}%
\definecolor{currentstroke}{rgb}{1.000000,1.000000,1.000000}%
\pgfsetstrokecolor{currentstroke}%
\pgfsetdash{}{0pt}%
\pgfpathmoveto{\pgfqpoint{1.872000in}{2.128000in}}%
\pgfpathlineto{\pgfqpoint{2.208000in}{2.128000in}}%
\pgfpathlineto{\pgfqpoint{2.208000in}{1.896000in}}%
\pgfpathlineto{\pgfqpoint{1.872000in}{1.896000in}}%
\pgfpathlineto{\pgfqpoint{1.872000in}{2.128000in}}%
\pgfusepath{fill}%
\end{pgfscope}%
\begin{pgfscope}%
\pgfpathrectangle{\pgfqpoint{1.200000in}{1.200000in}}{\pgfqpoint{3.360000in}{2.320000in}}%
\pgfusepath{clip}%
\pgfsetbuttcap%
\pgfsetroundjoin%
\definecolor{currentfill}{rgb}{0.355676,0.119109,0.316210}%
\pgfsetfillcolor{currentfill}%
\pgfsetlinewidth{0.000000pt}%
\definecolor{currentstroke}{rgb}{1.000000,1.000000,1.000000}%
\pgfsetstrokecolor{currentstroke}%
\pgfsetdash{}{0pt}%
\pgfpathmoveto{\pgfqpoint{2.208000in}{2.128000in}}%
\pgfpathlineto{\pgfqpoint{2.544000in}{2.128000in}}%
\pgfpathlineto{\pgfqpoint{2.544000in}{1.896000in}}%
\pgfpathlineto{\pgfqpoint{2.208000in}{1.896000in}}%
\pgfpathlineto{\pgfqpoint{2.208000in}{2.128000in}}%
\pgfusepath{fill}%
\end{pgfscope}%
\begin{pgfscope}%
\pgfpathrectangle{\pgfqpoint{1.200000in}{1.200000in}}{\pgfqpoint{3.360000in}{2.320000in}}%
\pgfusepath{clip}%
\pgfsetbuttcap%
\pgfsetroundjoin%
\definecolor{currentfill}{rgb}{0.312084,0.115053,0.298245}%
\pgfsetfillcolor{currentfill}%
\pgfsetlinewidth{0.000000pt}%
\definecolor{currentstroke}{rgb}{1.000000,1.000000,1.000000}%
\pgfsetstrokecolor{currentstroke}%
\pgfsetdash{}{0pt}%
\pgfpathmoveto{\pgfqpoint{2.544000in}{2.128000in}}%
\pgfpathlineto{\pgfqpoint{2.880000in}{2.128000in}}%
\pgfpathlineto{\pgfqpoint{2.880000in}{1.896000in}}%
\pgfpathlineto{\pgfqpoint{2.544000in}{1.896000in}}%
\pgfpathlineto{\pgfqpoint{2.544000in}{2.128000in}}%
\pgfusepath{fill}%
\end{pgfscope}%
\begin{pgfscope}%
\pgfpathrectangle{\pgfqpoint{1.200000in}{1.200000in}}{\pgfqpoint{3.360000in}{2.320000in}}%
\pgfusepath{clip}%
\pgfsetbuttcap%
\pgfsetroundjoin%
\definecolor{currentfill}{rgb}{0.324459,0.116396,0.303694}%
\pgfsetfillcolor{currentfill}%
\pgfsetlinewidth{0.000000pt}%
\definecolor{currentstroke}{rgb}{1.000000,1.000000,1.000000}%
\pgfsetstrokecolor{currentstroke}%
\pgfsetdash{}{0pt}%
\pgfpathmoveto{\pgfqpoint{2.880000in}{2.128000in}}%
\pgfpathlineto{\pgfqpoint{3.216000in}{2.128000in}}%
\pgfpathlineto{\pgfqpoint{3.216000in}{1.896000in}}%
\pgfpathlineto{\pgfqpoint{2.880000in}{1.896000in}}%
\pgfpathlineto{\pgfqpoint{2.880000in}{2.128000in}}%
\pgfusepath{fill}%
\end{pgfscope}%
\begin{pgfscope}%
\pgfpathrectangle{\pgfqpoint{1.200000in}{1.200000in}}{\pgfqpoint{3.360000in}{2.320000in}}%
\pgfusepath{clip}%
\pgfsetbuttcap%
\pgfsetroundjoin%
\definecolor{currentfill}{rgb}{0.981377,0.920617,0.865369}%
\pgfsetfillcolor{currentfill}%
\pgfsetlinewidth{0.000000pt}%
\definecolor{currentstroke}{rgb}{1.000000,1.000000,1.000000}%
\pgfsetstrokecolor{currentstroke}%
\pgfsetdash{}{0pt}%
\pgfpathmoveto{\pgfqpoint{3.216000in}{2.128000in}}%
\pgfpathlineto{\pgfqpoint{3.552000in}{2.128000in}}%
\pgfpathlineto{\pgfqpoint{3.552000in}{1.896000in}}%
\pgfpathlineto{\pgfqpoint{3.216000in}{1.896000in}}%
\pgfpathlineto{\pgfqpoint{3.216000in}{2.128000in}}%
\pgfusepath{fill}%
\end{pgfscope}%
\begin{pgfscope}%
\pgfpathrectangle{\pgfqpoint{1.200000in}{1.200000in}}{\pgfqpoint{3.360000in}{2.320000in}}%
\pgfusepath{clip}%
\pgfsetbuttcap%
\pgfsetroundjoin%
\definecolor{currentfill}{rgb}{0.355676,0.119109,0.316210}%
\pgfsetfillcolor{currentfill}%
\pgfsetlinewidth{0.000000pt}%
\definecolor{currentstroke}{rgb}{1.000000,1.000000,1.000000}%
\pgfsetstrokecolor{currentstroke}%
\pgfsetdash{}{0pt}%
\pgfpathmoveto{\pgfqpoint{3.552000in}{2.128000in}}%
\pgfpathlineto{\pgfqpoint{3.888000in}{2.128000in}}%
\pgfpathlineto{\pgfqpoint{3.888000in}{1.896000in}}%
\pgfpathlineto{\pgfqpoint{3.552000in}{1.896000in}}%
\pgfpathlineto{\pgfqpoint{3.552000in}{2.128000in}}%
\pgfusepath{fill}%
\end{pgfscope}%
\begin{pgfscope}%
\pgfpathrectangle{\pgfqpoint{1.200000in}{1.200000in}}{\pgfqpoint{3.360000in}{2.320000in}}%
\pgfusepath{clip}%
\pgfsetbuttcap%
\pgfsetroundjoin%
\definecolor{currentfill}{rgb}{0.368272,0.119937,0.320767}%
\pgfsetfillcolor{currentfill}%
\pgfsetlinewidth{0.000000pt}%
\definecolor{currentstroke}{rgb}{1.000000,1.000000,1.000000}%
\pgfsetstrokecolor{currentstroke}%
\pgfsetdash{}{0pt}%
\pgfpathmoveto{\pgfqpoint{3.888000in}{2.128000in}}%
\pgfpathlineto{\pgfqpoint{4.224000in}{2.128000in}}%
\pgfpathlineto{\pgfqpoint{4.224000in}{1.896000in}}%
\pgfpathlineto{\pgfqpoint{3.888000in}{1.896000in}}%
\pgfpathlineto{\pgfqpoint{3.888000in}{2.128000in}}%
\pgfusepath{fill}%
\end{pgfscope}%
\begin{pgfscope}%
\pgfpathrectangle{\pgfqpoint{1.200000in}{1.200000in}}{\pgfqpoint{3.360000in}{2.320000in}}%
\pgfusepath{clip}%
\pgfsetbuttcap%
\pgfsetroundjoin%
\definecolor{currentfill}{rgb}{0.368272,0.119937,0.320767}%
\pgfsetfillcolor{currentfill}%
\pgfsetlinewidth{0.000000pt}%
\definecolor{currentstroke}{rgb}{1.000000,1.000000,1.000000}%
\pgfsetstrokecolor{currentstroke}%
\pgfsetdash{}{0pt}%
\pgfpathmoveto{\pgfqpoint{4.224000in}{2.128000in}}%
\pgfpathlineto{\pgfqpoint{4.560000in}{2.128000in}}%
\pgfpathlineto{\pgfqpoint{4.560000in}{1.896000in}}%
\pgfpathlineto{\pgfqpoint{4.224000in}{1.896000in}}%
\pgfpathlineto{\pgfqpoint{4.224000in}{2.128000in}}%
\pgfusepath{fill}%
\end{pgfscope}%
\begin{pgfscope}%
\pgfpathrectangle{\pgfqpoint{1.200000in}{1.200000in}}{\pgfqpoint{3.360000in}{2.320000in}}%
\pgfusepath{clip}%
\pgfsetbuttcap%
\pgfsetroundjoin%
\definecolor{currentfill}{rgb}{0.438608,0.121770,0.341195}%
\pgfsetfillcolor{currentfill}%
\pgfsetlinewidth{0.000000pt}%
\definecolor{currentstroke}{rgb}{1.000000,1.000000,1.000000}%
\pgfsetstrokecolor{currentstroke}%
\pgfsetdash{}{0pt}%
\pgfpathmoveto{\pgfqpoint{1.200000in}{1.896000in}}%
\pgfpathlineto{\pgfqpoint{1.536000in}{1.896000in}}%
\pgfpathlineto{\pgfqpoint{1.536000in}{1.664000in}}%
\pgfpathlineto{\pgfqpoint{1.200000in}{1.664000in}}%
\pgfpathlineto{\pgfqpoint{1.200000in}{1.896000in}}%
\pgfusepath{fill}%
\end{pgfscope}%
\begin{pgfscope}%
\pgfpathrectangle{\pgfqpoint{1.200000in}{1.200000in}}{\pgfqpoint{3.360000in}{2.320000in}}%
\pgfusepath{clip}%
\pgfsetbuttcap%
\pgfsetroundjoin%
\definecolor{currentfill}{rgb}{0.361966,0.119541,0.318520}%
\pgfsetfillcolor{currentfill}%
\pgfsetlinewidth{0.000000pt}%
\definecolor{currentstroke}{rgb}{1.000000,1.000000,1.000000}%
\pgfsetstrokecolor{currentstroke}%
\pgfsetdash{}{0pt}%
\pgfpathmoveto{\pgfqpoint{1.536000in}{1.896000in}}%
\pgfpathlineto{\pgfqpoint{1.872000in}{1.896000in}}%
\pgfpathlineto{\pgfqpoint{1.872000in}{1.664000in}}%
\pgfpathlineto{\pgfqpoint{1.536000in}{1.664000in}}%
\pgfpathlineto{\pgfqpoint{1.536000in}{1.896000in}}%
\pgfusepath{fill}%
\end{pgfscope}%
\begin{pgfscope}%
\pgfpathrectangle{\pgfqpoint{1.200000in}{1.200000in}}{\pgfqpoint{3.360000in}{2.320000in}}%
\pgfusepath{clip}%
\pgfsetbuttcap%
\pgfsetroundjoin%
\definecolor{currentfill}{rgb}{0.355676,0.119109,0.316210}%
\pgfsetfillcolor{currentfill}%
\pgfsetlinewidth{0.000000pt}%
\definecolor{currentstroke}{rgb}{1.000000,1.000000,1.000000}%
\pgfsetstrokecolor{currentstroke}%
\pgfsetdash{}{0pt}%
\pgfpathmoveto{\pgfqpoint{1.872000in}{1.896000in}}%
\pgfpathlineto{\pgfqpoint{2.208000in}{1.896000in}}%
\pgfpathlineto{\pgfqpoint{2.208000in}{1.664000in}}%
\pgfpathlineto{\pgfqpoint{1.872000in}{1.664000in}}%
\pgfpathlineto{\pgfqpoint{1.872000in}{1.896000in}}%
\pgfusepath{fill}%
\end{pgfscope}%
\begin{pgfscope}%
\pgfpathrectangle{\pgfqpoint{1.200000in}{1.200000in}}{\pgfqpoint{3.360000in}{2.320000in}}%
\pgfusepath{clip}%
\pgfsetbuttcap%
\pgfsetroundjoin%
\definecolor{currentfill}{rgb}{0.343142,0.118134,0.311397}%
\pgfsetfillcolor{currentfill}%
\pgfsetlinewidth{0.000000pt}%
\definecolor{currentstroke}{rgb}{1.000000,1.000000,1.000000}%
\pgfsetstrokecolor{currentstroke}%
\pgfsetdash{}{0pt}%
\pgfpathmoveto{\pgfqpoint{2.208000in}{1.896000in}}%
\pgfpathlineto{\pgfqpoint{2.544000in}{1.896000in}}%
\pgfpathlineto{\pgfqpoint{2.544000in}{1.664000in}}%
\pgfpathlineto{\pgfqpoint{2.208000in}{1.664000in}}%
\pgfpathlineto{\pgfqpoint{2.208000in}{1.896000in}}%
\pgfusepath{fill}%
\end{pgfscope}%
\begin{pgfscope}%
\pgfpathrectangle{\pgfqpoint{1.200000in}{1.200000in}}{\pgfqpoint{3.360000in}{2.320000in}}%
\pgfusepath{clip}%
\pgfsetbuttcap%
\pgfsetroundjoin%
\definecolor{currentfill}{rgb}{0.380929,0.120615,0.325065}%
\pgfsetfillcolor{currentfill}%
\pgfsetlinewidth{0.000000pt}%
\definecolor{currentstroke}{rgb}{1.000000,1.000000,1.000000}%
\pgfsetstrokecolor{currentstroke}%
\pgfsetdash{}{0pt}%
\pgfpathmoveto{\pgfqpoint{2.544000in}{1.896000in}}%
\pgfpathlineto{\pgfqpoint{2.880000in}{1.896000in}}%
\pgfpathlineto{\pgfqpoint{2.880000in}{1.664000in}}%
\pgfpathlineto{\pgfqpoint{2.544000in}{1.664000in}}%
\pgfpathlineto{\pgfqpoint{2.544000in}{1.896000in}}%
\pgfusepath{fill}%
\end{pgfscope}%
\begin{pgfscope}%
\pgfpathrectangle{\pgfqpoint{1.200000in}{1.200000in}}{\pgfqpoint{3.360000in}{2.320000in}}%
\pgfusepath{clip}%
\pgfsetbuttcap%
\pgfsetroundjoin%
\definecolor{currentfill}{rgb}{0.355676,0.119109,0.316210}%
\pgfsetfillcolor{currentfill}%
\pgfsetlinewidth{0.000000pt}%
\definecolor{currentstroke}{rgb}{1.000000,1.000000,1.000000}%
\pgfsetstrokecolor{currentstroke}%
\pgfsetdash{}{0pt}%
\pgfpathmoveto{\pgfqpoint{2.880000in}{1.896000in}}%
\pgfpathlineto{\pgfqpoint{3.216000in}{1.896000in}}%
\pgfpathlineto{\pgfqpoint{3.216000in}{1.664000in}}%
\pgfpathlineto{\pgfqpoint{2.880000in}{1.664000in}}%
\pgfpathlineto{\pgfqpoint{2.880000in}{1.896000in}}%
\pgfusepath{fill}%
\end{pgfscope}%
\begin{pgfscope}%
\pgfpathrectangle{\pgfqpoint{1.200000in}{1.200000in}}{\pgfqpoint{3.360000in}{2.320000in}}%
\pgfusepath{clip}%
\pgfsetbuttcap%
\pgfsetroundjoin%
\definecolor{currentfill}{rgb}{0.355676,0.119109,0.316210}%
\pgfsetfillcolor{currentfill}%
\pgfsetlinewidth{0.000000pt}%
\definecolor{currentstroke}{rgb}{1.000000,1.000000,1.000000}%
\pgfsetstrokecolor{currentstroke}%
\pgfsetdash{}{0pt}%
\pgfpathmoveto{\pgfqpoint{3.216000in}{1.896000in}}%
\pgfpathlineto{\pgfqpoint{3.552000in}{1.896000in}}%
\pgfpathlineto{\pgfqpoint{3.552000in}{1.664000in}}%
\pgfpathlineto{\pgfqpoint{3.216000in}{1.664000in}}%
\pgfpathlineto{\pgfqpoint{3.216000in}{1.896000in}}%
\pgfusepath{fill}%
\end{pgfscope}%
\begin{pgfscope}%
\pgfpathrectangle{\pgfqpoint{1.200000in}{1.200000in}}{\pgfqpoint{3.360000in}{2.320000in}}%
\pgfusepath{clip}%
\pgfsetbuttcap%
\pgfsetroundjoin%
\definecolor{currentfill}{rgb}{0.981377,0.920617,0.865369}%
\pgfsetfillcolor{currentfill}%
\pgfsetlinewidth{0.000000pt}%
\definecolor{currentstroke}{rgb}{1.000000,1.000000,1.000000}%
\pgfsetstrokecolor{currentstroke}%
\pgfsetdash{}{0pt}%
\pgfpathmoveto{\pgfqpoint{3.552000in}{1.896000in}}%
\pgfpathlineto{\pgfqpoint{3.888000in}{1.896000in}}%
\pgfpathlineto{\pgfqpoint{3.888000in}{1.664000in}}%
\pgfpathlineto{\pgfqpoint{3.552000in}{1.664000in}}%
\pgfpathlineto{\pgfqpoint{3.552000in}{1.896000in}}%
\pgfusepath{fill}%
\end{pgfscope}%
\begin{pgfscope}%
\pgfpathrectangle{\pgfqpoint{1.200000in}{1.200000in}}{\pgfqpoint{3.360000in}{2.320000in}}%
\pgfusepath{clip}%
\pgfsetbuttcap%
\pgfsetroundjoin%
\definecolor{currentfill}{rgb}{0.349401,0.118640,0.313835}%
\pgfsetfillcolor{currentfill}%
\pgfsetlinewidth{0.000000pt}%
\definecolor{currentstroke}{rgb}{1.000000,1.000000,1.000000}%
\pgfsetstrokecolor{currentstroke}%
\pgfsetdash{}{0pt}%
\pgfpathmoveto{\pgfqpoint{3.888000in}{1.896000in}}%
\pgfpathlineto{\pgfqpoint{4.224000in}{1.896000in}}%
\pgfpathlineto{\pgfqpoint{4.224000in}{1.664000in}}%
\pgfpathlineto{\pgfqpoint{3.888000in}{1.664000in}}%
\pgfpathlineto{\pgfqpoint{3.888000in}{1.896000in}}%
\pgfusepath{fill}%
\end{pgfscope}%
\begin{pgfscope}%
\pgfpathrectangle{\pgfqpoint{1.200000in}{1.200000in}}{\pgfqpoint{3.360000in}{2.320000in}}%
\pgfusepath{clip}%
\pgfsetbuttcap%
\pgfsetroundjoin%
\definecolor{currentfill}{rgb}{0.361966,0.119541,0.318520}%
\pgfsetfillcolor{currentfill}%
\pgfsetlinewidth{0.000000pt}%
\definecolor{currentstroke}{rgb}{1.000000,1.000000,1.000000}%
\pgfsetstrokecolor{currentstroke}%
\pgfsetdash{}{0pt}%
\pgfpathmoveto{\pgfqpoint{4.224000in}{1.896000in}}%
\pgfpathlineto{\pgfqpoint{4.560000in}{1.896000in}}%
\pgfpathlineto{\pgfqpoint{4.560000in}{1.664000in}}%
\pgfpathlineto{\pgfqpoint{4.224000in}{1.664000in}}%
\pgfpathlineto{\pgfqpoint{4.224000in}{1.896000in}}%
\pgfusepath{fill}%
\end{pgfscope}%
\begin{pgfscope}%
\pgfpathrectangle{\pgfqpoint{1.200000in}{1.200000in}}{\pgfqpoint{3.360000in}{2.320000in}}%
\pgfusepath{clip}%
\pgfsetbuttcap%
\pgfsetroundjoin%
\definecolor{currentfill}{rgb}{0.368272,0.119937,0.320767}%
\pgfsetfillcolor{currentfill}%
\pgfsetlinewidth{0.000000pt}%
\definecolor{currentstroke}{rgb}{1.000000,1.000000,1.000000}%
\pgfsetstrokecolor{currentstroke}%
\pgfsetdash{}{0pt}%
\pgfpathmoveto{\pgfqpoint{1.200000in}{1.664000in}}%
\pgfpathlineto{\pgfqpoint{1.536000in}{1.664000in}}%
\pgfpathlineto{\pgfqpoint{1.536000in}{1.432000in}}%
\pgfpathlineto{\pgfqpoint{1.200000in}{1.432000in}}%
\pgfpathlineto{\pgfqpoint{1.200000in}{1.664000in}}%
\pgfusepath{fill}%
\end{pgfscope}%
\begin{pgfscope}%
\pgfpathrectangle{\pgfqpoint{1.200000in}{1.200000in}}{\pgfqpoint{3.360000in}{2.320000in}}%
\pgfusepath{clip}%
\pgfsetbuttcap%
\pgfsetroundjoin%
\definecolor{currentfill}{rgb}{0.336898,0.117591,0.308894}%
\pgfsetfillcolor{currentfill}%
\pgfsetlinewidth{0.000000pt}%
\definecolor{currentstroke}{rgb}{1.000000,1.000000,1.000000}%
\pgfsetstrokecolor{currentstroke}%
\pgfsetdash{}{0pt}%
\pgfpathmoveto{\pgfqpoint{1.536000in}{1.664000in}}%
\pgfpathlineto{\pgfqpoint{1.872000in}{1.664000in}}%
\pgfpathlineto{\pgfqpoint{1.872000in}{1.432000in}}%
\pgfpathlineto{\pgfqpoint{1.536000in}{1.432000in}}%
\pgfpathlineto{\pgfqpoint{1.536000in}{1.664000in}}%
\pgfusepath{fill}%
\end{pgfscope}%
\begin{pgfscope}%
\pgfpathrectangle{\pgfqpoint{1.200000in}{1.200000in}}{\pgfqpoint{3.360000in}{2.320000in}}%
\pgfusepath{clip}%
\pgfsetbuttcap%
\pgfsetroundjoin%
\definecolor{currentfill}{rgb}{0.361966,0.119541,0.318520}%
\pgfsetfillcolor{currentfill}%
\pgfsetlinewidth{0.000000pt}%
\definecolor{currentstroke}{rgb}{1.000000,1.000000,1.000000}%
\pgfsetstrokecolor{currentstroke}%
\pgfsetdash{}{0pt}%
\pgfpathmoveto{\pgfqpoint{1.872000in}{1.664000in}}%
\pgfpathlineto{\pgfqpoint{2.208000in}{1.664000in}}%
\pgfpathlineto{\pgfqpoint{2.208000in}{1.432000in}}%
\pgfpathlineto{\pgfqpoint{1.872000in}{1.432000in}}%
\pgfpathlineto{\pgfqpoint{1.872000in}{1.664000in}}%
\pgfusepath{fill}%
\end{pgfscope}%
\begin{pgfscope}%
\pgfpathrectangle{\pgfqpoint{1.200000in}{1.200000in}}{\pgfqpoint{3.360000in}{2.320000in}}%
\pgfusepath{clip}%
\pgfsetbuttcap%
\pgfsetroundjoin%
\definecolor{currentfill}{rgb}{0.336898,0.117591,0.308894}%
\pgfsetfillcolor{currentfill}%
\pgfsetlinewidth{0.000000pt}%
\definecolor{currentstroke}{rgb}{1.000000,1.000000,1.000000}%
\pgfsetstrokecolor{currentstroke}%
\pgfsetdash{}{0pt}%
\pgfpathmoveto{\pgfqpoint{2.208000in}{1.664000in}}%
\pgfpathlineto{\pgfqpoint{2.544000in}{1.664000in}}%
\pgfpathlineto{\pgfqpoint{2.544000in}{1.432000in}}%
\pgfpathlineto{\pgfqpoint{2.208000in}{1.432000in}}%
\pgfpathlineto{\pgfqpoint{2.208000in}{1.664000in}}%
\pgfusepath{fill}%
\end{pgfscope}%
\begin{pgfscope}%
\pgfpathrectangle{\pgfqpoint{1.200000in}{1.200000in}}{\pgfqpoint{3.360000in}{2.320000in}}%
\pgfusepath{clip}%
\pgfsetbuttcap%
\pgfsetroundjoin%
\definecolor{currentfill}{rgb}{0.343142,0.118134,0.311397}%
\pgfsetfillcolor{currentfill}%
\pgfsetlinewidth{0.000000pt}%
\definecolor{currentstroke}{rgb}{1.000000,1.000000,1.000000}%
\pgfsetstrokecolor{currentstroke}%
\pgfsetdash{}{0pt}%
\pgfpathmoveto{\pgfqpoint{2.544000in}{1.664000in}}%
\pgfpathlineto{\pgfqpoint{2.880000in}{1.664000in}}%
\pgfpathlineto{\pgfqpoint{2.880000in}{1.432000in}}%
\pgfpathlineto{\pgfqpoint{2.544000in}{1.432000in}}%
\pgfpathlineto{\pgfqpoint{2.544000in}{1.664000in}}%
\pgfusepath{fill}%
\end{pgfscope}%
\begin{pgfscope}%
\pgfpathrectangle{\pgfqpoint{1.200000in}{1.200000in}}{\pgfqpoint{3.360000in}{2.320000in}}%
\pgfusepath{clip}%
\pgfsetbuttcap%
\pgfsetroundjoin%
\definecolor{currentfill}{rgb}{0.343142,0.118134,0.311397}%
\pgfsetfillcolor{currentfill}%
\pgfsetlinewidth{0.000000pt}%
\definecolor{currentstroke}{rgb}{1.000000,1.000000,1.000000}%
\pgfsetstrokecolor{currentstroke}%
\pgfsetdash{}{0pt}%
\pgfpathmoveto{\pgfqpoint{2.880000in}{1.664000in}}%
\pgfpathlineto{\pgfqpoint{3.216000in}{1.664000in}}%
\pgfpathlineto{\pgfqpoint{3.216000in}{1.432000in}}%
\pgfpathlineto{\pgfqpoint{2.880000in}{1.432000in}}%
\pgfpathlineto{\pgfqpoint{2.880000in}{1.664000in}}%
\pgfusepath{fill}%
\end{pgfscope}%
\begin{pgfscope}%
\pgfpathrectangle{\pgfqpoint{1.200000in}{1.200000in}}{\pgfqpoint{3.360000in}{2.320000in}}%
\pgfusepath{clip}%
\pgfsetbuttcap%
\pgfsetroundjoin%
\definecolor{currentfill}{rgb}{0.368272,0.119937,0.320767}%
\pgfsetfillcolor{currentfill}%
\pgfsetlinewidth{0.000000pt}%
\definecolor{currentstroke}{rgb}{1.000000,1.000000,1.000000}%
\pgfsetstrokecolor{currentstroke}%
\pgfsetdash{}{0pt}%
\pgfpathmoveto{\pgfqpoint{3.216000in}{1.664000in}}%
\pgfpathlineto{\pgfqpoint{3.552000in}{1.664000in}}%
\pgfpathlineto{\pgfqpoint{3.552000in}{1.432000in}}%
\pgfpathlineto{\pgfqpoint{3.216000in}{1.432000in}}%
\pgfpathlineto{\pgfqpoint{3.216000in}{1.664000in}}%
\pgfusepath{fill}%
\end{pgfscope}%
\begin{pgfscope}%
\pgfpathrectangle{\pgfqpoint{1.200000in}{1.200000in}}{\pgfqpoint{3.360000in}{2.320000in}}%
\pgfusepath{clip}%
\pgfsetbuttcap%
\pgfsetroundjoin%
\definecolor{currentfill}{rgb}{0.349401,0.118640,0.313835}%
\pgfsetfillcolor{currentfill}%
\pgfsetlinewidth{0.000000pt}%
\definecolor{currentstroke}{rgb}{1.000000,1.000000,1.000000}%
\pgfsetstrokecolor{currentstroke}%
\pgfsetdash{}{0pt}%
\pgfpathmoveto{\pgfqpoint{3.552000in}{1.664000in}}%
\pgfpathlineto{\pgfqpoint{3.888000in}{1.664000in}}%
\pgfpathlineto{\pgfqpoint{3.888000in}{1.432000in}}%
\pgfpathlineto{\pgfqpoint{3.552000in}{1.432000in}}%
\pgfpathlineto{\pgfqpoint{3.552000in}{1.664000in}}%
\pgfusepath{fill}%
\end{pgfscope}%
\begin{pgfscope}%
\pgfpathrectangle{\pgfqpoint{1.200000in}{1.200000in}}{\pgfqpoint{3.360000in}{2.320000in}}%
\pgfusepath{clip}%
\pgfsetbuttcap%
\pgfsetroundjoin%
\definecolor{currentfill}{rgb}{0.981377,0.920617,0.865369}%
\pgfsetfillcolor{currentfill}%
\pgfsetlinewidth{0.000000pt}%
\definecolor{currentstroke}{rgb}{1.000000,1.000000,1.000000}%
\pgfsetstrokecolor{currentstroke}%
\pgfsetdash{}{0pt}%
\pgfpathmoveto{\pgfqpoint{3.888000in}{1.664000in}}%
\pgfpathlineto{\pgfqpoint{4.224000in}{1.664000in}}%
\pgfpathlineto{\pgfqpoint{4.224000in}{1.432000in}}%
\pgfpathlineto{\pgfqpoint{3.888000in}{1.432000in}}%
\pgfpathlineto{\pgfqpoint{3.888000in}{1.664000in}}%
\pgfusepath{fill}%
\end{pgfscope}%
\begin{pgfscope}%
\pgfpathrectangle{\pgfqpoint{1.200000in}{1.200000in}}{\pgfqpoint{3.360000in}{2.320000in}}%
\pgfusepath{clip}%
\pgfsetbuttcap%
\pgfsetroundjoin%
\definecolor{currentfill}{rgb}{0.349401,0.118640,0.313835}%
\pgfsetfillcolor{currentfill}%
\pgfsetlinewidth{0.000000pt}%
\definecolor{currentstroke}{rgb}{1.000000,1.000000,1.000000}%
\pgfsetstrokecolor{currentstroke}%
\pgfsetdash{}{0pt}%
\pgfpathmoveto{\pgfqpoint{4.224000in}{1.664000in}}%
\pgfpathlineto{\pgfqpoint{4.560000in}{1.664000in}}%
\pgfpathlineto{\pgfqpoint{4.560000in}{1.432000in}}%
\pgfpathlineto{\pgfqpoint{4.224000in}{1.432000in}}%
\pgfpathlineto{\pgfqpoint{4.224000in}{1.664000in}}%
\pgfusepath{fill}%
\end{pgfscope}%
\begin{pgfscope}%
\pgfpathrectangle{\pgfqpoint{1.200000in}{1.200000in}}{\pgfqpoint{3.360000in}{2.320000in}}%
\pgfusepath{clip}%
\pgfsetbuttcap%
\pgfsetroundjoin%
\definecolor{currentfill}{rgb}{0.355676,0.119109,0.316210}%
\pgfsetfillcolor{currentfill}%
\pgfsetlinewidth{0.000000pt}%
\definecolor{currentstroke}{rgb}{1.000000,1.000000,1.000000}%
\pgfsetstrokecolor{currentstroke}%
\pgfsetdash{}{0pt}%
\pgfpathmoveto{\pgfqpoint{1.200000in}{1.432000in}}%
\pgfpathlineto{\pgfqpoint{1.536000in}{1.432000in}}%
\pgfpathlineto{\pgfqpoint{1.536000in}{1.200000in}}%
\pgfpathlineto{\pgfqpoint{1.200000in}{1.200000in}}%
\pgfpathlineto{\pgfqpoint{1.200000in}{1.432000in}}%
\pgfusepath{fill}%
\end{pgfscope}%
\begin{pgfscope}%
\pgfpathrectangle{\pgfqpoint{1.200000in}{1.200000in}}{\pgfqpoint{3.360000in}{2.320000in}}%
\pgfusepath{clip}%
\pgfsetbuttcap%
\pgfsetroundjoin%
\definecolor{currentfill}{rgb}{0.380929,0.120615,0.325065}%
\pgfsetfillcolor{currentfill}%
\pgfsetlinewidth{0.000000pt}%
\definecolor{currentstroke}{rgb}{1.000000,1.000000,1.000000}%
\pgfsetstrokecolor{currentstroke}%
\pgfsetdash{}{0pt}%
\pgfpathmoveto{\pgfqpoint{1.536000in}{1.432000in}}%
\pgfpathlineto{\pgfqpoint{1.872000in}{1.432000in}}%
\pgfpathlineto{\pgfqpoint{1.872000in}{1.200000in}}%
\pgfpathlineto{\pgfqpoint{1.536000in}{1.200000in}}%
\pgfpathlineto{\pgfqpoint{1.536000in}{1.432000in}}%
\pgfusepath{fill}%
\end{pgfscope}%
\begin{pgfscope}%
\pgfpathrectangle{\pgfqpoint{1.200000in}{1.200000in}}{\pgfqpoint{3.360000in}{2.320000in}}%
\pgfusepath{clip}%
\pgfsetbuttcap%
\pgfsetroundjoin%
\definecolor{currentfill}{rgb}{0.349401,0.118640,0.313835}%
\pgfsetfillcolor{currentfill}%
\pgfsetlinewidth{0.000000pt}%
\definecolor{currentstroke}{rgb}{1.000000,1.000000,1.000000}%
\pgfsetstrokecolor{currentstroke}%
\pgfsetdash{}{0pt}%
\pgfpathmoveto{\pgfqpoint{1.872000in}{1.432000in}}%
\pgfpathlineto{\pgfqpoint{2.208000in}{1.432000in}}%
\pgfpathlineto{\pgfqpoint{2.208000in}{1.200000in}}%
\pgfpathlineto{\pgfqpoint{1.872000in}{1.200000in}}%
\pgfpathlineto{\pgfqpoint{1.872000in}{1.432000in}}%
\pgfusepath{fill}%
\end{pgfscope}%
\begin{pgfscope}%
\pgfpathrectangle{\pgfqpoint{1.200000in}{1.200000in}}{\pgfqpoint{3.360000in}{2.320000in}}%
\pgfusepath{clip}%
\pgfsetbuttcap%
\pgfsetroundjoin%
\definecolor{currentfill}{rgb}{0.343142,0.118134,0.311397}%
\pgfsetfillcolor{currentfill}%
\pgfsetlinewidth{0.000000pt}%
\definecolor{currentstroke}{rgb}{1.000000,1.000000,1.000000}%
\pgfsetstrokecolor{currentstroke}%
\pgfsetdash{}{0pt}%
\pgfpathmoveto{\pgfqpoint{2.208000in}{1.432000in}}%
\pgfpathlineto{\pgfqpoint{2.544000in}{1.432000in}}%
\pgfpathlineto{\pgfqpoint{2.544000in}{1.200000in}}%
\pgfpathlineto{\pgfqpoint{2.208000in}{1.200000in}}%
\pgfpathlineto{\pgfqpoint{2.208000in}{1.432000in}}%
\pgfusepath{fill}%
\end{pgfscope}%
\begin{pgfscope}%
\pgfpathrectangle{\pgfqpoint{1.200000in}{1.200000in}}{\pgfqpoint{3.360000in}{2.320000in}}%
\pgfusepath{clip}%
\pgfsetbuttcap%
\pgfsetroundjoin%
\definecolor{currentfill}{rgb}{0.343142,0.118134,0.311397}%
\pgfsetfillcolor{currentfill}%
\pgfsetlinewidth{0.000000pt}%
\definecolor{currentstroke}{rgb}{1.000000,1.000000,1.000000}%
\pgfsetstrokecolor{currentstroke}%
\pgfsetdash{}{0pt}%
\pgfpathmoveto{\pgfqpoint{2.544000in}{1.432000in}}%
\pgfpathlineto{\pgfqpoint{2.880000in}{1.432000in}}%
\pgfpathlineto{\pgfqpoint{2.880000in}{1.200000in}}%
\pgfpathlineto{\pgfqpoint{2.544000in}{1.200000in}}%
\pgfpathlineto{\pgfqpoint{2.544000in}{1.432000in}}%
\pgfusepath{fill}%
\end{pgfscope}%
\begin{pgfscope}%
\pgfpathrectangle{\pgfqpoint{1.200000in}{1.200000in}}{\pgfqpoint{3.360000in}{2.320000in}}%
\pgfusepath{clip}%
\pgfsetbuttcap%
\pgfsetroundjoin%
\definecolor{currentfill}{rgb}{0.349401,0.118640,0.313835}%
\pgfsetfillcolor{currentfill}%
\pgfsetlinewidth{0.000000pt}%
\definecolor{currentstroke}{rgb}{1.000000,1.000000,1.000000}%
\pgfsetstrokecolor{currentstroke}%
\pgfsetdash{}{0pt}%
\pgfpathmoveto{\pgfqpoint{2.880000in}{1.432000in}}%
\pgfpathlineto{\pgfqpoint{3.216000in}{1.432000in}}%
\pgfpathlineto{\pgfqpoint{3.216000in}{1.200000in}}%
\pgfpathlineto{\pgfqpoint{2.880000in}{1.200000in}}%
\pgfpathlineto{\pgfqpoint{2.880000in}{1.432000in}}%
\pgfusepath{fill}%
\end{pgfscope}%
\begin{pgfscope}%
\pgfpathrectangle{\pgfqpoint{1.200000in}{1.200000in}}{\pgfqpoint{3.360000in}{2.320000in}}%
\pgfusepath{clip}%
\pgfsetbuttcap%
\pgfsetroundjoin%
\definecolor{currentfill}{rgb}{0.368272,0.119937,0.320767}%
\pgfsetfillcolor{currentfill}%
\pgfsetlinewidth{0.000000pt}%
\definecolor{currentstroke}{rgb}{1.000000,1.000000,1.000000}%
\pgfsetstrokecolor{currentstroke}%
\pgfsetdash{}{0pt}%
\pgfpathmoveto{\pgfqpoint{3.216000in}{1.432000in}}%
\pgfpathlineto{\pgfqpoint{3.552000in}{1.432000in}}%
\pgfpathlineto{\pgfqpoint{3.552000in}{1.200000in}}%
\pgfpathlineto{\pgfqpoint{3.216000in}{1.200000in}}%
\pgfpathlineto{\pgfqpoint{3.216000in}{1.432000in}}%
\pgfusepath{fill}%
\end{pgfscope}%
\begin{pgfscope}%
\pgfpathrectangle{\pgfqpoint{1.200000in}{1.200000in}}{\pgfqpoint{3.360000in}{2.320000in}}%
\pgfusepath{clip}%
\pgfsetbuttcap%
\pgfsetroundjoin%
\definecolor{currentfill}{rgb}{0.361966,0.119541,0.318520}%
\pgfsetfillcolor{currentfill}%
\pgfsetlinewidth{0.000000pt}%
\definecolor{currentstroke}{rgb}{1.000000,1.000000,1.000000}%
\pgfsetstrokecolor{currentstroke}%
\pgfsetdash{}{0pt}%
\pgfpathmoveto{\pgfqpoint{3.552000in}{1.432000in}}%
\pgfpathlineto{\pgfqpoint{3.888000in}{1.432000in}}%
\pgfpathlineto{\pgfqpoint{3.888000in}{1.200000in}}%
\pgfpathlineto{\pgfqpoint{3.552000in}{1.200000in}}%
\pgfpathlineto{\pgfqpoint{3.552000in}{1.432000in}}%
\pgfusepath{fill}%
\end{pgfscope}%
\begin{pgfscope}%
\pgfpathrectangle{\pgfqpoint{1.200000in}{1.200000in}}{\pgfqpoint{3.360000in}{2.320000in}}%
\pgfusepath{clip}%
\pgfsetbuttcap%
\pgfsetroundjoin%
\definecolor{currentfill}{rgb}{0.349401,0.118640,0.313835}%
\pgfsetfillcolor{currentfill}%
\pgfsetlinewidth{0.000000pt}%
\definecolor{currentstroke}{rgb}{1.000000,1.000000,1.000000}%
\pgfsetstrokecolor{currentstroke}%
\pgfsetdash{}{0pt}%
\pgfpathmoveto{\pgfqpoint{3.888000in}{1.432000in}}%
\pgfpathlineto{\pgfqpoint{4.224000in}{1.432000in}}%
\pgfpathlineto{\pgfqpoint{4.224000in}{1.200000in}}%
\pgfpathlineto{\pgfqpoint{3.888000in}{1.200000in}}%
\pgfpathlineto{\pgfqpoint{3.888000in}{1.432000in}}%
\pgfusepath{fill}%
\end{pgfscope}%
\begin{pgfscope}%
\pgfpathrectangle{\pgfqpoint{1.200000in}{1.200000in}}{\pgfqpoint{3.360000in}{2.320000in}}%
\pgfusepath{clip}%
\pgfsetbuttcap%
\pgfsetroundjoin%
\definecolor{currentfill}{rgb}{0.981377,0.920617,0.865369}%
\pgfsetfillcolor{currentfill}%
\pgfsetlinewidth{0.000000pt}%
\definecolor{currentstroke}{rgb}{1.000000,1.000000,1.000000}%
\pgfsetstrokecolor{currentstroke}%
\pgfsetdash{}{0pt}%
\pgfpathmoveto{\pgfqpoint{4.224000in}{1.432000in}}%
\pgfpathlineto{\pgfqpoint{4.560000in}{1.432000in}}%
\pgfpathlineto{\pgfqpoint{4.560000in}{1.200000in}}%
\pgfpathlineto{\pgfqpoint{4.224000in}{1.200000in}}%
\pgfpathlineto{\pgfqpoint{4.224000in}{1.432000in}}%
\pgfusepath{fill}%
\end{pgfscope}%
\begin{pgfscope}%
\pgfsetbuttcap%
\pgfsetroundjoin%
\definecolor{currentfill}{rgb}{0.000000,0.000000,0.000000}%
\pgfsetfillcolor{currentfill}%
\pgfsetlinewidth{0.803000pt}%
\definecolor{currentstroke}{rgb}{0.000000,0.000000,0.000000}%
\pgfsetstrokecolor{currentstroke}%
\pgfsetdash{}{0pt}%
\pgfsys@defobject{currentmarker}{\pgfqpoint{0.000000in}{-0.048611in}}{\pgfqpoint{0.000000in}{0.000000in}}{%
\pgfpathmoveto{\pgfqpoint{0.000000in}{0.000000in}}%
\pgfpathlineto{\pgfqpoint{0.000000in}{-0.048611in}}%
\pgfusepath{stroke,fill}%
}%
\begin{pgfscope}%
\pgfsys@transformshift{1.368000in}{1.200000in}%
\pgfsys@useobject{currentmarker}{}%
\end{pgfscope}%
\end{pgfscope}%
\begin{pgfscope}%
\definecolor{textcolor}{rgb}{0.000000,0.000000,0.000000}%
\pgfsetstrokecolor{textcolor}%
\pgfsetfillcolor{textcolor}%
\pgftext[x=1.402722in, y=0.626389in, left, base,rotate=90.000000]{\color{textcolor}\rmfamily\fontsize{10.000000}{12.000000}\selectfont Balance}%
\end{pgfscope}%
\begin{pgfscope}%
\pgfsetbuttcap%
\pgfsetroundjoin%
\definecolor{currentfill}{rgb}{0.000000,0.000000,0.000000}%
\pgfsetfillcolor{currentfill}%
\pgfsetlinewidth{0.803000pt}%
\definecolor{currentstroke}{rgb}{0.000000,0.000000,0.000000}%
\pgfsetstrokecolor{currentstroke}%
\pgfsetdash{}{0pt}%
\pgfsys@defobject{currentmarker}{\pgfqpoint{0.000000in}{-0.048611in}}{\pgfqpoint{0.000000in}{0.000000in}}{%
\pgfpathmoveto{\pgfqpoint{0.000000in}{0.000000in}}%
\pgfpathlineto{\pgfqpoint{0.000000in}{-0.048611in}}%
\pgfusepath{stroke,fill}%
}%
\begin{pgfscope}%
\pgfsys@transformshift{1.704000in}{1.200000in}%
\pgfsys@useobject{currentmarker}{}%
\end{pgfscope}%
\end{pgfscope}%
\begin{pgfscope}%
\definecolor{textcolor}{rgb}{0.000000,0.000000,0.000000}%
\pgfsetstrokecolor{textcolor}%
\pgfsetfillcolor{textcolor}%
\pgftext[x=1.738722in, y=0.114167in, left, base,rotate=90.000000]{\color{textcolor}\rmfamily\fontsize{10.000000}{12.000000}\selectfont IsActiveMember}%
\end{pgfscope}%
\begin{pgfscope}%
\pgfsetbuttcap%
\pgfsetroundjoin%
\definecolor{currentfill}{rgb}{0.000000,0.000000,0.000000}%
\pgfsetfillcolor{currentfill}%
\pgfsetlinewidth{0.803000pt}%
\definecolor{currentstroke}{rgb}{0.000000,0.000000,0.000000}%
\pgfsetstrokecolor{currentstroke}%
\pgfsetdash{}{0pt}%
\pgfsys@defobject{currentmarker}{\pgfqpoint{0.000000in}{-0.048611in}}{\pgfqpoint{0.000000in}{0.000000in}}{%
\pgfpathmoveto{\pgfqpoint{0.000000in}{0.000000in}}%
\pgfpathlineto{\pgfqpoint{0.000000in}{-0.048611in}}%
\pgfusepath{stroke,fill}%
}%
\begin{pgfscope}%
\pgfsys@transformshift{2.040000in}{1.200000in}%
\pgfsys@useobject{currentmarker}{}%
\end{pgfscope}%
\end{pgfscope}%
\begin{pgfscope}%
\definecolor{textcolor}{rgb}{0.000000,0.000000,0.000000}%
\pgfsetstrokecolor{textcolor}%
\pgfsetfillcolor{textcolor}%
\pgftext[x=2.074722in, y=0.685694in, left, base,rotate=90.000000]{\color{textcolor}\rmfamily\fontsize{10.000000}{12.000000}\selectfont Tenure}%
\end{pgfscope}%
\begin{pgfscope}%
\pgfsetbuttcap%
\pgfsetroundjoin%
\definecolor{currentfill}{rgb}{0.000000,0.000000,0.000000}%
\pgfsetfillcolor{currentfill}%
\pgfsetlinewidth{0.803000pt}%
\definecolor{currentstroke}{rgb}{0.000000,0.000000,0.000000}%
\pgfsetstrokecolor{currentstroke}%
\pgfsetdash{}{0pt}%
\pgfsys@defobject{currentmarker}{\pgfqpoint{0.000000in}{-0.048611in}}{\pgfqpoint{0.000000in}{0.000000in}}{%
\pgfpathmoveto{\pgfqpoint{0.000000in}{0.000000in}}%
\pgfpathlineto{\pgfqpoint{0.000000in}{-0.048611in}}%
\pgfusepath{stroke,fill}%
}%
\begin{pgfscope}%
\pgfsys@transformshift{2.376000in}{1.200000in}%
\pgfsys@useobject{currentmarker}{}%
\end{pgfscope}%
\end{pgfscope}%
\begin{pgfscope}%
\definecolor{textcolor}{rgb}{0.000000,0.000000,0.000000}%
\pgfsetstrokecolor{textcolor}%
\pgfsetfillcolor{textcolor}%
\pgftext[x=2.410722in, y=0.418333in, left, base,rotate=90.000000]{\color{textcolor}\rmfamily\fontsize{10.000000}{12.000000}\selectfont HasCrCard}%
\end{pgfscope}%
\begin{pgfscope}%
\pgfsetbuttcap%
\pgfsetroundjoin%
\definecolor{currentfill}{rgb}{0.000000,0.000000,0.000000}%
\pgfsetfillcolor{currentfill}%
\pgfsetlinewidth{0.803000pt}%
\definecolor{currentstroke}{rgb}{0.000000,0.000000,0.000000}%
\pgfsetstrokecolor{currentstroke}%
\pgfsetdash{}{0pt}%
\pgfsys@defobject{currentmarker}{\pgfqpoint{0.000000in}{-0.048611in}}{\pgfqpoint{0.000000in}{0.000000in}}{%
\pgfpathmoveto{\pgfqpoint{0.000000in}{0.000000in}}%
\pgfpathlineto{\pgfqpoint{0.000000in}{-0.048611in}}%
\pgfusepath{stroke,fill}%
}%
\begin{pgfscope}%
\pgfsys@transformshift{2.712000in}{1.200000in}%
\pgfsys@useobject{currentmarker}{}%
\end{pgfscope}%
\end{pgfscope}%
\begin{pgfscope}%
\definecolor{textcolor}{rgb}{0.000000,0.000000,0.000000}%
\pgfsetstrokecolor{textcolor}%
\pgfsetfillcolor{textcolor}%
\pgftext[x=2.747417in, y=0.867500in, left, base,rotate=90.000000]{\color{textcolor}\rmfamily\fontsize{10.000000}{12.000000}\selectfont Age}%
\end{pgfscope}%
\begin{pgfscope}%
\pgfsetbuttcap%
\pgfsetroundjoin%
\definecolor{currentfill}{rgb}{0.000000,0.000000,0.000000}%
\pgfsetfillcolor{currentfill}%
\pgfsetlinewidth{0.803000pt}%
\definecolor{currentstroke}{rgb}{0.000000,0.000000,0.000000}%
\pgfsetstrokecolor{currentstroke}%
\pgfsetdash{}{0pt}%
\pgfsys@defobject{currentmarker}{\pgfqpoint{0.000000in}{-0.048611in}}{\pgfqpoint{0.000000in}{0.000000in}}{%
\pgfpathmoveto{\pgfqpoint{0.000000in}{0.000000in}}%
\pgfpathlineto{\pgfqpoint{0.000000in}{-0.048611in}}%
\pgfusepath{stroke,fill}%
}%
\begin{pgfscope}%
\pgfsys@transformshift{3.048000in}{1.200000in}%
\pgfsys@useobject{currentmarker}{}%
\end{pgfscope}%
\end{pgfscope}%
\begin{pgfscope}%
\definecolor{textcolor}{rgb}{0.000000,0.000000,0.000000}%
\pgfsetstrokecolor{textcolor}%
\pgfsetfillcolor{textcolor}%
\pgftext[x=3.082722in, y=0.661528in, left, base,rotate=90.000000]{\color{textcolor}\rmfamily\fontsize{10.000000}{12.000000}\selectfont Gender}%
\end{pgfscope}%
\begin{pgfscope}%
\pgfsetbuttcap%
\pgfsetroundjoin%
\definecolor{currentfill}{rgb}{0.000000,0.000000,0.000000}%
\pgfsetfillcolor{currentfill}%
\pgfsetlinewidth{0.803000pt}%
\definecolor{currentstroke}{rgb}{0.000000,0.000000,0.000000}%
\pgfsetstrokecolor{currentstroke}%
\pgfsetdash{}{0pt}%
\pgfsys@defobject{currentmarker}{\pgfqpoint{0.000000in}{-0.048611in}}{\pgfqpoint{0.000000in}{0.000000in}}{%
\pgfpathmoveto{\pgfqpoint{0.000000in}{0.000000in}}%
\pgfpathlineto{\pgfqpoint{0.000000in}{-0.048611in}}%
\pgfusepath{stroke,fill}%
}%
\begin{pgfscope}%
\pgfsys@transformshift{3.384000in}{1.200000in}%
\pgfsys@useobject{currentmarker}{}%
\end{pgfscope}%
\end{pgfscope}%
\begin{pgfscope}%
\definecolor{textcolor}{rgb}{0.000000,0.000000,0.000000}%
\pgfsetstrokecolor{textcolor}%
\pgfsetfillcolor{textcolor}%
\pgftext[x=3.418722in, y=0.107917in, left, base,rotate=90.000000]{\color{textcolor}\rmfamily\fontsize{10.000000}{12.000000}\selectfont NumOfProducts}%
\end{pgfscope}%
\begin{pgfscope}%
\pgfsetbuttcap%
\pgfsetroundjoin%
\definecolor{currentfill}{rgb}{0.000000,0.000000,0.000000}%
\pgfsetfillcolor{currentfill}%
\pgfsetlinewidth{0.803000pt}%
\definecolor{currentstroke}{rgb}{0.000000,0.000000,0.000000}%
\pgfsetstrokecolor{currentstroke}%
\pgfsetdash{}{0pt}%
\pgfsys@defobject{currentmarker}{\pgfqpoint{0.000000in}{-0.048611in}}{\pgfqpoint{0.000000in}{0.000000in}}{%
\pgfpathmoveto{\pgfqpoint{0.000000in}{0.000000in}}%
\pgfpathlineto{\pgfqpoint{0.000000in}{-0.048611in}}%
\pgfusepath{stroke,fill}%
}%
\begin{pgfscope}%
\pgfsys@transformshift{3.720000in}{1.200000in}%
\pgfsys@useobject{currentmarker}{}%
\end{pgfscope}%
\end{pgfscope}%
\begin{pgfscope}%
\definecolor{textcolor}{rgb}{0.000000,0.000000,0.000000}%
\pgfsetstrokecolor{textcolor}%
\pgfsetfillcolor{textcolor}%
\pgftext[x=3.754653in, y=0.445417in, left, base,rotate=90.000000]{\color{textcolor}\rmfamily\fontsize{10.000000}{12.000000}\selectfont Geography}%
\end{pgfscope}%
\begin{pgfscope}%
\pgfsetbuttcap%
\pgfsetroundjoin%
\definecolor{currentfill}{rgb}{0.000000,0.000000,0.000000}%
\pgfsetfillcolor{currentfill}%
\pgfsetlinewidth{0.803000pt}%
\definecolor{currentstroke}{rgb}{0.000000,0.000000,0.000000}%
\pgfsetstrokecolor{currentstroke}%
\pgfsetdash{}{0pt}%
\pgfsys@defobject{currentmarker}{\pgfqpoint{0.000000in}{-0.048611in}}{\pgfqpoint{0.000000in}{0.000000in}}{%
\pgfpathmoveto{\pgfqpoint{0.000000in}{0.000000in}}%
\pgfpathlineto{\pgfqpoint{0.000000in}{-0.048611in}}%
\pgfusepath{stroke,fill}%
}%
\begin{pgfscope}%
\pgfsys@transformshift{4.056000in}{1.200000in}%
\pgfsys@useobject{currentmarker}{}%
\end{pgfscope}%
\end{pgfscope}%
\begin{pgfscope}%
\definecolor{textcolor}{rgb}{0.000000,0.000000,0.000000}%
\pgfsetstrokecolor{textcolor}%
\pgfsetfillcolor{textcolor}%
\pgftext[x=4.090722in, y=0.100278in, left, base,rotate=90.000000]{\color{textcolor}\rmfamily\fontsize{10.000000}{12.000000}\selectfont EstimatedSalary}%
\end{pgfscope}%
\begin{pgfscope}%
\pgfsetbuttcap%
\pgfsetroundjoin%
\definecolor{currentfill}{rgb}{0.000000,0.000000,0.000000}%
\pgfsetfillcolor{currentfill}%
\pgfsetlinewidth{0.803000pt}%
\definecolor{currentstroke}{rgb}{0.000000,0.000000,0.000000}%
\pgfsetstrokecolor{currentstroke}%
\pgfsetdash{}{0pt}%
\pgfsys@defobject{currentmarker}{\pgfqpoint{0.000000in}{-0.048611in}}{\pgfqpoint{0.000000in}{0.000000in}}{%
\pgfpathmoveto{\pgfqpoint{0.000000in}{0.000000in}}%
\pgfpathlineto{\pgfqpoint{0.000000in}{-0.048611in}}%
\pgfusepath{stroke,fill}%
}%
\begin{pgfscope}%
\pgfsys@transformshift{4.392000in}{1.200000in}%
\pgfsys@useobject{currentmarker}{}%
\end{pgfscope}%
\end{pgfscope}%
\begin{pgfscope}%
\definecolor{textcolor}{rgb}{0.000000,0.000000,0.000000}%
\pgfsetstrokecolor{textcolor}%
\pgfsetfillcolor{textcolor}%
\pgftext[x=4.426722in, y=0.392083in, left, base,rotate=90.000000]{\color{textcolor}\rmfamily\fontsize{10.000000}{12.000000}\selectfont CreditScore}%
\end{pgfscope}%
\begin{pgfscope}%
\pgfsetbuttcap%
\pgfsetroundjoin%
\definecolor{currentfill}{rgb}{0.000000,0.000000,0.000000}%
\pgfsetfillcolor{currentfill}%
\pgfsetlinewidth{0.803000pt}%
\definecolor{currentstroke}{rgb}{0.000000,0.000000,0.000000}%
\pgfsetstrokecolor{currentstroke}%
\pgfsetdash{}{0pt}%
\pgfsys@defobject{currentmarker}{\pgfqpoint{-0.048611in}{0.000000in}}{\pgfqpoint{-0.000000in}{0.000000in}}{%
\pgfpathmoveto{\pgfqpoint{-0.000000in}{0.000000in}}%
\pgfpathlineto{\pgfqpoint{-0.048611in}{0.000000in}}%
\pgfusepath{stroke,fill}%
}%
\begin{pgfscope}%
\pgfsys@transformshift{1.200000in}{3.404000in}%
\pgfsys@useobject{currentmarker}{}%
\end{pgfscope}%
\end{pgfscope}%
\begin{pgfscope}%
\definecolor{textcolor}{rgb}{0.000000,0.000000,0.000000}%
\pgfsetstrokecolor{textcolor}%
\pgfsetfillcolor{textcolor}%
\pgftext[x=1.102778in,y=3.404000in,right,]{\color{textcolor}\rmfamily\fontsize{10.000000}{12.000000}\selectfont Balance}%
\end{pgfscope}%
\begin{pgfscope}%
\pgfsetbuttcap%
\pgfsetroundjoin%
\definecolor{currentfill}{rgb}{0.000000,0.000000,0.000000}%
\pgfsetfillcolor{currentfill}%
\pgfsetlinewidth{0.803000pt}%
\definecolor{currentstroke}{rgb}{0.000000,0.000000,0.000000}%
\pgfsetstrokecolor{currentstroke}%
\pgfsetdash{}{0pt}%
\pgfsys@defobject{currentmarker}{\pgfqpoint{-0.048611in}{0.000000in}}{\pgfqpoint{-0.000000in}{0.000000in}}{%
\pgfpathmoveto{\pgfqpoint{-0.000000in}{0.000000in}}%
\pgfpathlineto{\pgfqpoint{-0.048611in}{0.000000in}}%
\pgfusepath{stroke,fill}%
}%
\begin{pgfscope}%
\pgfsys@transformshift{1.200000in}{3.172000in}%
\pgfsys@useobject{currentmarker}{}%
\end{pgfscope}%
\end{pgfscope}%
\begin{pgfscope}%
\definecolor{textcolor}{rgb}{0.000000,0.000000,0.000000}%
\pgfsetstrokecolor{textcolor}%
\pgfsetfillcolor{textcolor}%
\pgftext[x=1.102778in,y=3.172000in,right,]{\color{textcolor}\rmfamily\fontsize{10.000000}{12.000000}\selectfont IsActiveMember}%
\end{pgfscope}%
\begin{pgfscope}%
\pgfsetbuttcap%
\pgfsetroundjoin%
\definecolor{currentfill}{rgb}{0.000000,0.000000,0.000000}%
\pgfsetfillcolor{currentfill}%
\pgfsetlinewidth{0.803000pt}%
\definecolor{currentstroke}{rgb}{0.000000,0.000000,0.000000}%
\pgfsetstrokecolor{currentstroke}%
\pgfsetdash{}{0pt}%
\pgfsys@defobject{currentmarker}{\pgfqpoint{-0.048611in}{0.000000in}}{\pgfqpoint{-0.000000in}{0.000000in}}{%
\pgfpathmoveto{\pgfqpoint{-0.000000in}{0.000000in}}%
\pgfpathlineto{\pgfqpoint{-0.048611in}{0.000000in}}%
\pgfusepath{stroke,fill}%
}%
\begin{pgfscope}%
\pgfsys@transformshift{1.200000in}{2.940000in}%
\pgfsys@useobject{currentmarker}{}%
\end{pgfscope}%
\end{pgfscope}%
\begin{pgfscope}%
\definecolor{textcolor}{rgb}{0.000000,0.000000,0.000000}%
\pgfsetstrokecolor{textcolor}%
\pgfsetfillcolor{textcolor}%
\pgftext[x=1.102778in,y=2.940000in,right,]{\color{textcolor}\rmfamily\fontsize{10.000000}{12.000000}\selectfont Tenure}%
\end{pgfscope}%
\begin{pgfscope}%
\pgfsetbuttcap%
\pgfsetroundjoin%
\definecolor{currentfill}{rgb}{0.000000,0.000000,0.000000}%
\pgfsetfillcolor{currentfill}%
\pgfsetlinewidth{0.803000pt}%
\definecolor{currentstroke}{rgb}{0.000000,0.000000,0.000000}%
\pgfsetstrokecolor{currentstroke}%
\pgfsetdash{}{0pt}%
\pgfsys@defobject{currentmarker}{\pgfqpoint{-0.048611in}{0.000000in}}{\pgfqpoint{-0.000000in}{0.000000in}}{%
\pgfpathmoveto{\pgfqpoint{-0.000000in}{0.000000in}}%
\pgfpathlineto{\pgfqpoint{-0.048611in}{0.000000in}}%
\pgfusepath{stroke,fill}%
}%
\begin{pgfscope}%
\pgfsys@transformshift{1.200000in}{2.708000in}%
\pgfsys@useobject{currentmarker}{}%
\end{pgfscope}%
\end{pgfscope}%
\begin{pgfscope}%
\definecolor{textcolor}{rgb}{0.000000,0.000000,0.000000}%
\pgfsetstrokecolor{textcolor}%
\pgfsetfillcolor{textcolor}%
\pgftext[x=1.102778in,y=2.708000in,right,]{\color{textcolor}\rmfamily\fontsize{10.000000}{12.000000}\selectfont HasCrCard}%
\end{pgfscope}%
\begin{pgfscope}%
\pgfsetbuttcap%
\pgfsetroundjoin%
\definecolor{currentfill}{rgb}{0.000000,0.000000,0.000000}%
\pgfsetfillcolor{currentfill}%
\pgfsetlinewidth{0.803000pt}%
\definecolor{currentstroke}{rgb}{0.000000,0.000000,0.000000}%
\pgfsetstrokecolor{currentstroke}%
\pgfsetdash{}{0pt}%
\pgfsys@defobject{currentmarker}{\pgfqpoint{-0.048611in}{0.000000in}}{\pgfqpoint{-0.000000in}{0.000000in}}{%
\pgfpathmoveto{\pgfqpoint{-0.000000in}{0.000000in}}%
\pgfpathlineto{\pgfqpoint{-0.048611in}{0.000000in}}%
\pgfusepath{stroke,fill}%
}%
\begin{pgfscope}%
\pgfsys@transformshift{1.200000in}{2.476000in}%
\pgfsys@useobject{currentmarker}{}%
\end{pgfscope}%
\end{pgfscope}%
\begin{pgfscope}%
\definecolor{textcolor}{rgb}{0.000000,0.000000,0.000000}%
\pgfsetstrokecolor{textcolor}%
\pgfsetfillcolor{textcolor}%
\pgftext[x=1.102778in,y=2.476000in,right,]{\color{textcolor}\rmfamily\fontsize{10.000000}{12.000000}\selectfont Age}%
\end{pgfscope}%
\begin{pgfscope}%
\pgfsetbuttcap%
\pgfsetroundjoin%
\definecolor{currentfill}{rgb}{0.000000,0.000000,0.000000}%
\pgfsetfillcolor{currentfill}%
\pgfsetlinewidth{0.803000pt}%
\definecolor{currentstroke}{rgb}{0.000000,0.000000,0.000000}%
\pgfsetstrokecolor{currentstroke}%
\pgfsetdash{}{0pt}%
\pgfsys@defobject{currentmarker}{\pgfqpoint{-0.048611in}{0.000000in}}{\pgfqpoint{-0.000000in}{0.000000in}}{%
\pgfpathmoveto{\pgfqpoint{-0.000000in}{0.000000in}}%
\pgfpathlineto{\pgfqpoint{-0.048611in}{0.000000in}}%
\pgfusepath{stroke,fill}%
}%
\begin{pgfscope}%
\pgfsys@transformshift{1.200000in}{2.244000in}%
\pgfsys@useobject{currentmarker}{}%
\end{pgfscope}%
\end{pgfscope}%
\begin{pgfscope}%
\definecolor{textcolor}{rgb}{0.000000,0.000000,0.000000}%
\pgfsetstrokecolor{textcolor}%
\pgfsetfillcolor{textcolor}%
\pgftext[x=1.102778in,y=2.244000in,right,]{\color{textcolor}\rmfamily\fontsize{10.000000}{12.000000}\selectfont Gender}%
\end{pgfscope}%
\begin{pgfscope}%
\pgfsetbuttcap%
\pgfsetroundjoin%
\definecolor{currentfill}{rgb}{0.000000,0.000000,0.000000}%
\pgfsetfillcolor{currentfill}%
\pgfsetlinewidth{0.803000pt}%
\definecolor{currentstroke}{rgb}{0.000000,0.000000,0.000000}%
\pgfsetstrokecolor{currentstroke}%
\pgfsetdash{}{0pt}%
\pgfsys@defobject{currentmarker}{\pgfqpoint{-0.048611in}{0.000000in}}{\pgfqpoint{-0.000000in}{0.000000in}}{%
\pgfpathmoveto{\pgfqpoint{-0.000000in}{0.000000in}}%
\pgfpathlineto{\pgfqpoint{-0.048611in}{0.000000in}}%
\pgfusepath{stroke,fill}%
}%
\begin{pgfscope}%
\pgfsys@transformshift{1.200000in}{2.012000in}%
\pgfsys@useobject{currentmarker}{}%
\end{pgfscope}%
\end{pgfscope}%
\begin{pgfscope}%
\definecolor{textcolor}{rgb}{0.000000,0.000000,0.000000}%
\pgfsetstrokecolor{textcolor}%
\pgfsetfillcolor{textcolor}%
\pgftext[x=1.102778in,y=2.012000in,right,]{\color{textcolor}\rmfamily\fontsize{10.000000}{12.000000}\selectfont NumOfProducts}%
\end{pgfscope}%
\begin{pgfscope}%
\pgfsetbuttcap%
\pgfsetroundjoin%
\definecolor{currentfill}{rgb}{0.000000,0.000000,0.000000}%
\pgfsetfillcolor{currentfill}%
\pgfsetlinewidth{0.803000pt}%
\definecolor{currentstroke}{rgb}{0.000000,0.000000,0.000000}%
\pgfsetstrokecolor{currentstroke}%
\pgfsetdash{}{0pt}%
\pgfsys@defobject{currentmarker}{\pgfqpoint{-0.048611in}{0.000000in}}{\pgfqpoint{-0.000000in}{0.000000in}}{%
\pgfpathmoveto{\pgfqpoint{-0.000000in}{0.000000in}}%
\pgfpathlineto{\pgfqpoint{-0.048611in}{0.000000in}}%
\pgfusepath{stroke,fill}%
}%
\begin{pgfscope}%
\pgfsys@transformshift{1.200000in}{1.780000in}%
\pgfsys@useobject{currentmarker}{}%
\end{pgfscope}%
\end{pgfscope}%
\begin{pgfscope}%
\definecolor{textcolor}{rgb}{0.000000,0.000000,0.000000}%
\pgfsetstrokecolor{textcolor}%
\pgfsetfillcolor{textcolor}%
\pgftext[x=1.102778in,y=1.780000in,right,]{\color{textcolor}\rmfamily\fontsize{10.000000}{12.000000}\selectfont Geography}%
\end{pgfscope}%
\begin{pgfscope}%
\pgfsetbuttcap%
\pgfsetroundjoin%
\definecolor{currentfill}{rgb}{0.000000,0.000000,0.000000}%
\pgfsetfillcolor{currentfill}%
\pgfsetlinewidth{0.803000pt}%
\definecolor{currentstroke}{rgb}{0.000000,0.000000,0.000000}%
\pgfsetstrokecolor{currentstroke}%
\pgfsetdash{}{0pt}%
\pgfsys@defobject{currentmarker}{\pgfqpoint{-0.048611in}{0.000000in}}{\pgfqpoint{-0.000000in}{0.000000in}}{%
\pgfpathmoveto{\pgfqpoint{-0.000000in}{0.000000in}}%
\pgfpathlineto{\pgfqpoint{-0.048611in}{0.000000in}}%
\pgfusepath{stroke,fill}%
}%
\begin{pgfscope}%
\pgfsys@transformshift{1.200000in}{1.548000in}%
\pgfsys@useobject{currentmarker}{}%
\end{pgfscope}%
\end{pgfscope}%
\begin{pgfscope}%
\definecolor{textcolor}{rgb}{0.000000,0.000000,0.000000}%
\pgfsetstrokecolor{textcolor}%
\pgfsetfillcolor{textcolor}%
\pgftext[x=1.102778in,y=1.548000in,right,]{\color{textcolor}\rmfamily\fontsize{10.000000}{12.000000}\selectfont EstimatedSalary}%
\end{pgfscope}%
\begin{pgfscope}%
\pgfsetbuttcap%
\pgfsetroundjoin%
\definecolor{currentfill}{rgb}{0.000000,0.000000,0.000000}%
\pgfsetfillcolor{currentfill}%
\pgfsetlinewidth{0.803000pt}%
\definecolor{currentstroke}{rgb}{0.000000,0.000000,0.000000}%
\pgfsetstrokecolor{currentstroke}%
\pgfsetdash{}{0pt}%
\pgfsys@defobject{currentmarker}{\pgfqpoint{-0.048611in}{0.000000in}}{\pgfqpoint{-0.000000in}{0.000000in}}{%
\pgfpathmoveto{\pgfqpoint{-0.000000in}{0.000000in}}%
\pgfpathlineto{\pgfqpoint{-0.048611in}{0.000000in}}%
\pgfusepath{stroke,fill}%
}%
\begin{pgfscope}%
\pgfsys@transformshift{1.200000in}{1.316000in}%
\pgfsys@useobject{currentmarker}{}%
\end{pgfscope}%
\end{pgfscope}%
\begin{pgfscope}%
\definecolor{textcolor}{rgb}{0.000000,0.000000,0.000000}%
\pgfsetstrokecolor{textcolor}%
\pgfsetfillcolor{textcolor}%
\pgftext[x=1.102778in,y=1.316000in,right,]{\color{textcolor}\rmfamily\fontsize{10.000000}{12.000000}\selectfont CreditScore}%
\end{pgfscope}%
\begin{pgfscope}%
\definecolor{textcolor}{rgb}{0.000000,0.000000,0.000000}%
\pgfsetstrokecolor{textcolor}%
\pgfsetfillcolor{textcolor}%
\pgftext[x=2.880000in,y=3.603333in,,base]{\color{textcolor}\rmfamily\fontsize{12.000000}{14.400000}\selectfont Correlation Matrix}%
\end{pgfscope}%
\begin{pgfscope}%
\pgfsetbuttcap%
\pgfsetmiterjoin%
\definecolor{currentfill}{rgb}{1.000000,1.000000,1.000000}%
\pgfsetfillcolor{currentfill}%
\pgfsetlinewidth{0.000000pt}%
\definecolor{currentstroke}{rgb}{0.000000,0.000000,0.000000}%
\pgfsetstrokecolor{currentstroke}%
\pgfsetstrokeopacity{0.000000}%
\pgfsetdash{}{0pt}%
\pgfpathmoveto{\pgfqpoint{4.770000in}{1.200000in}}%
\pgfpathlineto{\pgfqpoint{4.886000in}{1.200000in}}%
\pgfpathlineto{\pgfqpoint{4.886000in}{3.520000in}}%
\pgfpathlineto{\pgfqpoint{4.770000in}{3.520000in}}%
\pgfpathlineto{\pgfqpoint{4.770000in}{1.200000in}}%
\pgfpathclose%
\pgfusepath{fill}%
\end{pgfscope}%
\begin{pgfscope}%
\pgfpathrectangle{\pgfqpoint{4.770000in}{1.200000in}}{\pgfqpoint{0.116000in}{2.320000in}}%
\pgfusepath{clip}%
\pgfsetbuttcap%
\pgfsetmiterjoin%
\definecolor{currentfill}{rgb}{1.000000,1.000000,1.000000}%
\pgfsetfillcolor{currentfill}%
\pgfsetlinewidth{0.010037pt}%
\definecolor{currentstroke}{rgb}{1.000000,1.000000,1.000000}%
\pgfsetstrokecolor{currentstroke}%
\pgfsetdash{}{0pt}%
\pgfusepath{stroke,fill}%
\end{pgfscope}%
\begin{pgfscope}%
\pgfsys@transformshift{4.763889in}{1.194444in}%
\pgftext[left,bottom]{\includegraphics[interpolate=true,width=0.125000in,height=2.319444in]{corr-img0.png}}%
\end{pgfscope}%
\begin{pgfscope}%
\pgfsetbuttcap%
\pgfsetroundjoin%
\definecolor{currentfill}{rgb}{0.000000,0.000000,0.000000}%
\pgfsetfillcolor{currentfill}%
\pgfsetlinewidth{0.803000pt}%
\definecolor{currentstroke}{rgb}{0.000000,0.000000,0.000000}%
\pgfsetstrokecolor{currentstroke}%
\pgfsetdash{}{0pt}%
\pgfsys@defobject{currentmarker}{\pgfqpoint{0.000000in}{0.000000in}}{\pgfqpoint{0.048611in}{0.000000in}}{%
\pgfpathmoveto{\pgfqpoint{0.000000in}{0.000000in}}%
\pgfpathlineto{\pgfqpoint{0.048611in}{0.000000in}}%
\pgfusepath{stroke,fill}%
}%
\begin{pgfscope}%
\pgfsys@transformshift{4.886000in}{1.385325in}%
\pgfsys@useobject{currentmarker}{}%
\end{pgfscope}%
\end{pgfscope}%
\begin{pgfscope}%
\definecolor{textcolor}{rgb}{0.000000,0.000000,0.000000}%
\pgfsetstrokecolor{textcolor}%
\pgfsetfillcolor{textcolor}%
\pgftext[x=4.983222in, y=1.337130in, left, base]{\color{textcolor}\rmfamily\fontsize{10.000000}{12.000000}\selectfont \(\displaystyle {\ensuremath{-}0.2}\)}%
\end{pgfscope}%
\begin{pgfscope}%
\pgfsetbuttcap%
\pgfsetroundjoin%
\definecolor{currentfill}{rgb}{0.000000,0.000000,0.000000}%
\pgfsetfillcolor{currentfill}%
\pgfsetlinewidth{0.803000pt}%
\definecolor{currentstroke}{rgb}{0.000000,0.000000,0.000000}%
\pgfsetstrokecolor{currentstroke}%
\pgfsetdash{}{0pt}%
\pgfsys@defobject{currentmarker}{\pgfqpoint{0.000000in}{0.000000in}}{\pgfqpoint{0.048611in}{0.000000in}}{%
\pgfpathmoveto{\pgfqpoint{0.000000in}{0.000000in}}%
\pgfpathlineto{\pgfqpoint{0.048611in}{0.000000in}}%
\pgfusepath{stroke,fill}%
}%
\begin{pgfscope}%
\pgfsys@transformshift{4.886000in}{1.741104in}%
\pgfsys@useobject{currentmarker}{}%
\end{pgfscope}%
\end{pgfscope}%
\begin{pgfscope}%
\definecolor{textcolor}{rgb}{0.000000,0.000000,0.000000}%
\pgfsetstrokecolor{textcolor}%
\pgfsetfillcolor{textcolor}%
\pgftext[x=4.983222in, y=1.692910in, left, base]{\color{textcolor}\rmfamily\fontsize{10.000000}{12.000000}\selectfont \(\displaystyle {0.0}\)}%
\end{pgfscope}%
\begin{pgfscope}%
\pgfsetbuttcap%
\pgfsetroundjoin%
\definecolor{currentfill}{rgb}{0.000000,0.000000,0.000000}%
\pgfsetfillcolor{currentfill}%
\pgfsetlinewidth{0.803000pt}%
\definecolor{currentstroke}{rgb}{0.000000,0.000000,0.000000}%
\pgfsetstrokecolor{currentstroke}%
\pgfsetdash{}{0pt}%
\pgfsys@defobject{currentmarker}{\pgfqpoint{0.000000in}{0.000000in}}{\pgfqpoint{0.048611in}{0.000000in}}{%
\pgfpathmoveto{\pgfqpoint{0.000000in}{0.000000in}}%
\pgfpathlineto{\pgfqpoint{0.048611in}{0.000000in}}%
\pgfusepath{stroke,fill}%
}%
\begin{pgfscope}%
\pgfsys@transformshift{4.886000in}{2.096883in}%
\pgfsys@useobject{currentmarker}{}%
\end{pgfscope}%
\end{pgfscope}%
\begin{pgfscope}%
\definecolor{textcolor}{rgb}{0.000000,0.000000,0.000000}%
\pgfsetstrokecolor{textcolor}%
\pgfsetfillcolor{textcolor}%
\pgftext[x=4.983222in, y=2.048689in, left, base]{\color{textcolor}\rmfamily\fontsize{10.000000}{12.000000}\selectfont \(\displaystyle {0.2}\)}%
\end{pgfscope}%
\begin{pgfscope}%
\pgfsetbuttcap%
\pgfsetroundjoin%
\definecolor{currentfill}{rgb}{0.000000,0.000000,0.000000}%
\pgfsetfillcolor{currentfill}%
\pgfsetlinewidth{0.803000pt}%
\definecolor{currentstroke}{rgb}{0.000000,0.000000,0.000000}%
\pgfsetstrokecolor{currentstroke}%
\pgfsetdash{}{0pt}%
\pgfsys@defobject{currentmarker}{\pgfqpoint{0.000000in}{0.000000in}}{\pgfqpoint{0.048611in}{0.000000in}}{%
\pgfpathmoveto{\pgfqpoint{0.000000in}{0.000000in}}%
\pgfpathlineto{\pgfqpoint{0.048611in}{0.000000in}}%
\pgfusepath{stroke,fill}%
}%
\begin{pgfscope}%
\pgfsys@transformshift{4.886000in}{2.452662in}%
\pgfsys@useobject{currentmarker}{}%
\end{pgfscope}%
\end{pgfscope}%
\begin{pgfscope}%
\definecolor{textcolor}{rgb}{0.000000,0.000000,0.000000}%
\pgfsetstrokecolor{textcolor}%
\pgfsetfillcolor{textcolor}%
\pgftext[x=4.983222in, y=2.404468in, left, base]{\color{textcolor}\rmfamily\fontsize{10.000000}{12.000000}\selectfont \(\displaystyle {0.4}\)}%
\end{pgfscope}%
\begin{pgfscope}%
\pgfsetbuttcap%
\pgfsetroundjoin%
\definecolor{currentfill}{rgb}{0.000000,0.000000,0.000000}%
\pgfsetfillcolor{currentfill}%
\pgfsetlinewidth{0.803000pt}%
\definecolor{currentstroke}{rgb}{0.000000,0.000000,0.000000}%
\pgfsetstrokecolor{currentstroke}%
\pgfsetdash{}{0pt}%
\pgfsys@defobject{currentmarker}{\pgfqpoint{0.000000in}{0.000000in}}{\pgfqpoint{0.048611in}{0.000000in}}{%
\pgfpathmoveto{\pgfqpoint{0.000000in}{0.000000in}}%
\pgfpathlineto{\pgfqpoint{0.048611in}{0.000000in}}%
\pgfusepath{stroke,fill}%
}%
\begin{pgfscope}%
\pgfsys@transformshift{4.886000in}{2.808442in}%
\pgfsys@useobject{currentmarker}{}%
\end{pgfscope}%
\end{pgfscope}%
\begin{pgfscope}%
\definecolor{textcolor}{rgb}{0.000000,0.000000,0.000000}%
\pgfsetstrokecolor{textcolor}%
\pgfsetfillcolor{textcolor}%
\pgftext[x=4.983222in, y=2.760247in, left, base]{\color{textcolor}\rmfamily\fontsize{10.000000}{12.000000}\selectfont \(\displaystyle {0.6}\)}%
\end{pgfscope}%
\begin{pgfscope}%
\pgfsetbuttcap%
\pgfsetroundjoin%
\definecolor{currentfill}{rgb}{0.000000,0.000000,0.000000}%
\pgfsetfillcolor{currentfill}%
\pgfsetlinewidth{0.803000pt}%
\definecolor{currentstroke}{rgb}{0.000000,0.000000,0.000000}%
\pgfsetstrokecolor{currentstroke}%
\pgfsetdash{}{0pt}%
\pgfsys@defobject{currentmarker}{\pgfqpoint{0.000000in}{0.000000in}}{\pgfqpoint{0.048611in}{0.000000in}}{%
\pgfpathmoveto{\pgfqpoint{0.000000in}{0.000000in}}%
\pgfpathlineto{\pgfqpoint{0.048611in}{0.000000in}}%
\pgfusepath{stroke,fill}%
}%
\begin{pgfscope}%
\pgfsys@transformshift{4.886000in}{3.164221in}%
\pgfsys@useobject{currentmarker}{}%
\end{pgfscope}%
\end{pgfscope}%
\begin{pgfscope}%
\definecolor{textcolor}{rgb}{0.000000,0.000000,0.000000}%
\pgfsetstrokecolor{textcolor}%
\pgfsetfillcolor{textcolor}%
\pgftext[x=4.983222in, y=3.116026in, left, base]{\color{textcolor}\rmfamily\fontsize{10.000000}{12.000000}\selectfont \(\displaystyle {0.8}\)}%
\end{pgfscope}%
\begin{pgfscope}%
\pgfsetbuttcap%
\pgfsetroundjoin%
\definecolor{currentfill}{rgb}{0.000000,0.000000,0.000000}%
\pgfsetfillcolor{currentfill}%
\pgfsetlinewidth{0.803000pt}%
\definecolor{currentstroke}{rgb}{0.000000,0.000000,0.000000}%
\pgfsetstrokecolor{currentstroke}%
\pgfsetdash{}{0pt}%
\pgfsys@defobject{currentmarker}{\pgfqpoint{0.000000in}{0.000000in}}{\pgfqpoint{0.048611in}{0.000000in}}{%
\pgfpathmoveto{\pgfqpoint{0.000000in}{0.000000in}}%
\pgfpathlineto{\pgfqpoint{0.048611in}{0.000000in}}%
\pgfusepath{stroke,fill}%
}%
\begin{pgfscope}%
\pgfsys@transformshift{4.886000in}{3.520000in}%
\pgfsys@useobject{currentmarker}{}%
\end{pgfscope}%
\end{pgfscope}%
\begin{pgfscope}%
\definecolor{textcolor}{rgb}{0.000000,0.000000,0.000000}%
\pgfsetstrokecolor{textcolor}%
\pgfsetfillcolor{textcolor}%
\pgftext[x=4.983222in, y=3.471806in, left, base]{\color{textcolor}\rmfamily\fontsize{10.000000}{12.000000}\selectfont \(\displaystyle {1.0}\)}%
\end{pgfscope}%
\begin{pgfscope}%
\pgfsetrectcap%
\pgfsetmiterjoin%
\pgfsetlinewidth{0.000000pt}%
\definecolor{currentstroke}{rgb}{0.000000,0.000000,0.000000}%
\pgfsetstrokecolor{currentstroke}%
\pgfsetdash{}{0pt}%
\pgfpathmoveto{\pgfqpoint{4.770000in}{1.200000in}}%
\pgfpathlineto{\pgfqpoint{4.828000in}{1.200000in}}%
\pgfpathlineto{\pgfqpoint{4.886000in}{1.200000in}}%
\pgfpathlineto{\pgfqpoint{4.886000in}{3.520000in}}%
\pgfpathlineto{\pgfqpoint{4.828000in}{3.520000in}}%
\pgfpathlineto{\pgfqpoint{4.770000in}{3.520000in}}%
\pgfpathlineto{\pgfqpoint{4.770000in}{1.200000in}}%
\pgfpathclose%
\pgfusepath{}%
\end{pgfscope}%
\end{pgfpicture}%
\makeatother%
\endgroup%

    }
  \end{center}
  \caption{Correlations among predictor variables.}
\end{figure}

\section{Feature Selection}

\lettrine[nindent=0em,lines=2]{I} n order to remove excess “noise” 
from our model, we embark on feature selection. 
Specifically, we implement the feature selection process known as the Boruta algorithm. 
Boruta functions by creating a copy of each variable in the dataset, randomizing the copy, and attaching it back to the dataset. 
For example, if a dataset has original columns A and B, Boruta takes that dataset and produces a new dataset with columns \textit{A}, \textit{A\_Copy}, \textit{B}, and \textit{B\_Copy}, where \textit{A\_Copy} and \textit{B\_Copy} are shuffled versions of A and B, respectively (consider the copy columns essentially a random number generator sampling from the range of their original columns). 
Table \ref{table:boruta} illustrates this example.
\begin{table}
  \centering
  \caption{Boruta Example}
  \begin{tabular}{lllll}
  \toprule
   \textit{A} &  \textit{A\_copy} &  \textit{B} &  \textit{B\_copy} \\
  \midrule
   1 &       5 & 22 &      30 \\
   5 &       6 & 25 &      21 \\
   6 &       7 & 21 &      25 \\
   7 &       1 & 30 &      22 \\
  \bottomrule
  \end{tabular}
  \label{table:boruta}
  \end{table}
  
Next, Boruta takes all of the ``copy" columns and selects the copy column that was the best predictor of the dependent variable. 
If any of the original columns is a worse indicator of the dependent variable than the best copy column, 
that original column is removed. This entire process is repeated to ensure statistical significance. 
The foundational idea is that no variable in the data set is a worse predictor than a random column of data – each variable must have a non-random contribution towards predicting the dependent variable.

Ultimately, we whittle our original set of variables down to seven variables: age, estimated salary,  credit score, country of origin, number of products owned, member status, and remaining balance. We use these variables in the next step of our analysis, model implementation, tuning, and evaluation.

%------------------------------------------------

\section{Model Implementation, Tuning, \& Evaluation}
\lettrine[nindent=0em,lines=2]{T}o build an effective model using our dataset, we take a rapid prototyping approach by optimizing and comparing various supervised classification models in their ability to predict churn. Below we detail steps in our generalized approach to assessing each model, including train/test split, hyperparameter tuning, and model evaluation:

\noindent\subsection{Evaluation Metric} Considering the costly impacts of both false positives (wasting resources on customers not at risk of churn) and false negatives (losing customers), a reasonable metric that helps us balance between the two in model optimization is the F1 score. Specifically, we look at the F1 score for our “positive” predictions (overall ability to predict churn) as well as the macro-averaged F1 (unweighted average of F1 scores on predicting churn and retention, respectively) as a secondary metric.

\noindent\subsection{Train/Test Split} We used an 80/20 split for training/testing (holdout). The 20\% holdout set was used only to compute a final evaluation metric for a given model after optimizing on the training set.

\noindent\subsection{Standardization} For models that require standardization (e.g. Logistic Regression, SVM, KNN), we center the data and scale to unit variance.

\noindent \subsection{Hyperparameter Tuning} For each model, we identify its most important hyperparameters and define a reasonable search space for each parameter. We use grid search or randomized search to approximate an optimal combination of hyperparameters based on F1 score.

\noindent \subsection{Final Evaluation} For each optimized model, we compute its F1 score by comparing its predictions on the test set features against the true values of the test set target variable. F1 scores for five different models are reported in Table \ref{table:modelperf}

\section{Results}
\begin{table}
  \centering
  \caption{Model Performance}
  \begin{tabular}{lll}
  \toprule
                   Model &    F1 & Macro Averaged F1 \\
  \midrule
      Logistic Regresion &   TBD &               TBD \\
     K-Nearest Neighbors &  0.55 &              0.73 \\
  Support Vector Machine &  0.54 &              0.73 \\
           Random Forest &   0.6 &              0.75 \\
                 XGBoost &  0.57 &              0.75 \\
  \bottomrule
  \end{tabular}
  \label{table:modelperf}
  \end{table}
  
  
    

% \blindtext % Dummy text

% \begin{equation}
% \label{eq:emc}
% e = mc^2
% \end{equation}

% \blindtext % Dummy text

%------------------------------------------------

\section{Discussion}

\subsection{Subsection One}

% A statement requiring citation \cite{Figueredo:2009dg}.
% \blindtext % Dummy text

\subsection{Subsection Two}

% \blindtext % Dummy text

%----------------------------------------------------------------------------------------
%	REFERENCE LIST
%----------------------------------------------------------------------------------------

% \begin{thebibliography}{99} % Bibliography - this is intentionally simple in this template

% \bibitem[Figueredo and Wolf, 2009]{Figueredo:2009dg}
% Figueredo, A.~J. and Wolf, P. S.~A. (2009).
% \newblock Assortative pairing and life history strategy - a cross-cultural
%   study.
% \newblock {\em Human Nature}, 20:317--330.
 
% \end{thebibliography}

%----------------------------------------------------------------------------------------

\end{document}
